% =====================================================================
% main.tex -- HSC Physics ESSENTIALS, Module 5: Advanced Mechanics.
% Condensed, knowledge-only companion to the full Revision Guide
% (src/mod5/main.tex). Same document class and design system.
%
% Build with:
%     ./build.sh mod5_essentials
%
% Book order (Essentials editions):
%   cover - How to use these Essentials - Contents
%   chapters ch1.tex ... ch5.tex
%   All formulas at a glance                                      (formulas.tex)
%   Key experiments at a glance                                   (experiments.tex)
%   Glossary                                                      (glossary.tex)
% No questions, no answers, no practice set, no "Answering HSC questions" part.
% A missing part file is skipped, so a half-finished book still compiles.
% =====================================================================
\documentclass[11pt,a4paper,oneside,openany]{book}
\makeatletter\def\input@path{{./}{../common/}}\makeatother
\usepackage{hscphys}
\IfFileExists{./module_info.tex}{% module_info.tex -- Module 5 ESSENTIALS edition (condensed, knowledge only, no questions).
% Final PDF: build_book.sh mod5_essentials Module5_Advanced_Mechanics_Essentials.pdf  ->  output/Module5_Advanced_Mechanics_Essentials.pdf
\def\modulenumber{5}
\def\moduletitle{Advanced Mechanics}
\makeatletter
\renewcommand\hsccoversubtitle{HSC Physics Essentials — key ideas and formulas}
\renewcommand\hsccoverchips{%
  \hsc@coverchip{hscbigidea}{\faLightbulb}{Key ideas}\hspace{2mm}%
  \hsc@coverchip{hscformula}{\faIcon{square-root-alt}}{Formulas}\hspace{2mm}%
  \hsc@coverchip{hscexp}{\faFlask}{Key experiments}\hspace{2mm}%
  \hsc@coverchip{hscbigidea}{\faIcon{book}}{Glossary}}
\renewcommand\hsccoverfooter{Condensed edition \textperiodcentered\ knowledge points only \textperiodcentered\ no questions}
\makeatother
}{}
% PDF title of this edition (hscphys.sty sets the full guide's title; this runs after it)
\AtBeginDocument{\hypersetup{
  pdftitle={Module \modulenumber: \moduletitle\space - HSC Physics Essentials}}}

\begin{document}

%% ---------------------------------------------------------------- front
\makecover

\unnumberedchapter[\faInfoCircle]{How to use these Essentials}

This is the short, knowledge-only companion to the full \textbf{HSC Physics Revision Guide}
for Module~\modulenumber\ (\moduletitle). It covers every syllabus dot point of the NSW Physics
Stage~6 Syllabus (2017) as knowledge points only: key ideas, key words, facts, formulas and key
experiments. There are no questions here. The worked examples, Quick checks and the practice
set with full answers are in the full guide.

\section*{What the coloured boxes mean}

{\renewcommand\arraystretch{1.25}\arrayrulecolor{hscrule}
\begin{tabular}{@{}L{44mm}@{\hspace{4mm}}L{\dimexpr\linewidth-48mm\relax}@{}}
  \textcolor{hscbigidea}{\faLightbulb}\enspace\textbf{The big idea} (violet)
    & The main point of a topic in 1--3 sentences. Read it first.\\\hline
  \textcolor{hsckw}{\faKey}\enspace\textbf{Key words} (teal)
    & The most important terms, each with a one-line meaning. A new word is also printed in
      \term{bold} the first time it appears.\\\hline
  \textcolor{hscformula}{\faIcon{square-root-alt}}\enspace\textbf{Formula card} (blue)
    & The formula (exactly as on the official sheet, when it is printed there), each symbol
      with its SI unit, when to use it, where it comes from and a common trap.\\\hline
  \textcolor{hscexp}{\faFlask}\enspace\textbf{Key experiment} (brown)
    & Aim, what was done, the result, the conclusion and why it matters.\\\hline
  \textcolor{hscbrand}{\faClipboardList}\enspace\textbf{Summary} (navy)
    & ``In this chapter'' at the start of a chapter; ``Remember'' at the end: the facts to
      know by heart.\\\hline
\end{tabular}}

\smallskip
\noindent The badge under a formula tells you about the official formula sheet:
\begin{itemize}
  \item \hscbadgeonsheet\ \ it is printed on the sheet, so you are given it in the exam;
  \item \hscbadgenotonsheet\ \ learn it, or learn how to derive it;
  \item \hscbadgecombine{two formulas}\ \ join them.
\end{itemize}

\section*{How to study with these Essentials}

\begin{enumerate}
  \item Read \textbf{one chapter a day}: the big idea and the key words first, then the facts,
        the diagrams and the formula cards. Practise each chapter with the full guide.
  \item Use the tables at the back to see the whole module on a few pages: \emph{All formulas
        at a glance}, \emph{Key experiments at a glance} and the \emph{Glossary} (with Chinese).
  \item \textbf{The night before the exam} (Thursday 5 November 2026), re-read the
        ``Remember'' box at the end of each chapter.
\end{enumerate}

% Contents: smaller gaps between the entries than between text paragraphs, so that the
% contents of this short book fit on one page (local to this group; the style is unchanged).
{\setlength\parskip{0pt plus 1pt}\setstretch{1.15}\tableofcontents}

%% ---------------------------------------------------------------- chapters
\IfFileExists{./ch1.tex}{\setcounter{chapter}{0}% ch1.tex -- Module 5 ESSENTIALS edition, Chapter 1: Projectile Motion.
% (python3 -I). Convention in the whole chapter: origin at the launch point, up and forwards
% positive, so a_x = 0 and a_y = -g = -9.8 m s^-2.
\chapter{Projectile Motion}\label{ch1:chap}
\covers{M5.IQ1.01, M5.IQ1.02, M5.IQ1.03, M5.IQ1.04}

\begin{summarybox}{In this chapter}
\begin{itemize}
  \item The projectile model: two independent motions that share only the time.
  \item Maximum height, time of flight, range and final velocity.
  \item Energy in a flight, and what air resistance changes.
  \item Uniform velocity and constant acceleration models; the key experiment.
\end{itemize}
\end{summarybox}

%% ==================================================================
\section{The projectile model}\label{ch1:sec:model}
\covers{M5.IQ1.01, M5.IQ1.02, M5.IQ1.04}

\begin{bigidea}
A thrown ball moves sideways and up-and-down at the same time. The two motions do not affect
each other. Sideways the velocity stays the same; up and down, gravity changes it steadily.
\end{bigidea}

\begin{keywords}
  \kw{Projectile}{an object that moves through the air after it is thrown or launched, with no
    engine pushing it.}
  \kw{Component}{the part of a vector (an arrow quantity) that points in one direction.}
\end{keywords}

The projectile \term{model} (a simple description we can calculate with) has two
\term{assumptions}: \textbf{zero air resistance}, and a \textbf{constant} $g = \qty{9.8}{m.s^{-2}}$
straight down.

\begin{sidebyside}[0.42]
  \centering
  % scale 4.8 cm per metre in both directions; B thrown at 1.5 m/s; positions every 0.1 s
  \begin{tikzpicture}[x=0.8cm, y=0.8cm]
    \fill[ground] (-1.9,-5.08) rectangle (4.1,-4.88);
    \draw[hscink, line width=0.6pt] (-1.9,-4.88) -- (4.1,-4.88);
    \foreach \y in {0, -0.294, -1.176, -2.646, -4.704}
      \draw[guide] (-1.4,\y) -- (3.95,\y);
    \foreach \y/\t in {0/0, -1.176/0.2, -2.646/0.3, -4.704/0.4}
      \node[diagramlabel, anchor=east, font=\scriptsize, inner sep=1pt] at (-1.45,\y) {\t\,s};
    % the 0.1 s row is too close to the 0 s row for a label on the left: label it on the right
    \node[diagramlabel, anchor=west, font=\scriptsize, inner sep=1pt] at (4.0,-0.294) {0.1\,s};
    \foreach \y in {0, -0.294, -1.176, -2.646, -4.704}
      \fill[object] (-0.9,\y) circle (2.8pt);
    \foreach \x/\y in {0/0, 0.9/-0.294, 1.8/-1.176, 2.7/-2.646, 3.6/-4.704}
      \fill[object, fill=hscseriesB!35] (\x,\y) circle (2.8pt);
    \draw[velocity] (0.15,0.35) -- (0.9,0.35) node[right, text=hscvel, inner sep=1pt] {$u$};
    \node[diagramlabel, above, inner sep=2pt] at (-0.9,0.12) {\textbf{A}};
    \node[diagramlabel, above, inner sep=2pt] at (0,0.12) {\textbf{B}};
  \end{tikzpicture}
  \captionof{figure}{A is dropped. B is thrown sideways at the same moment.
    Positions every \qty{0.1}{s}.}\label{ch1:fig:strobe}
\nextside
  \begin{itemize}
    \item Gravity acts only downwards: \textbf{no horizontal force}. No horizontal force means
          no horizontal acceleration (Newton's second law), so \textbf{$v_x$ stays the same}.
    \item Vertically, $v_y$ changes by \qty{9.8}{m.s^{-1}} every second.
    \item The motions are \term{independent} (they do not affect each other) and share only
          the \emph{time} $t$. A and B in \cref{ch1:fig:strobe} land \textbf{at the same time}.
  \end{itemize}
\end{sidebyside}

\begin{itemize}
  \item The sideways distance grows with $t$; the drop grows with $t^2$. So the path (the
        \term{trajectory}) is a \term{parabola}: a curve symmetrical about its top.
  \item Signs: the \term{origin} is the launch point; up and forwards are positive. So
        $a_y = -g = \qty{-9.8}{m.s^{-2}}$ for the whole flight, even at the top.
\end{itemize}

\begin{formulacard}[notonsheet]{Components of the launch velocity}
  \cardformula{u_x = u\cos\theta \qquad u_y = u\sin\theta}
  \begin{symbols}
    \sym{$u$, $\theta$}{launch speed; launch angle up from the horizontal}{\si{m.s^{-1}};
      \si{\degree} (SI: rad)}
    \sym{$u_x$, $u_y$}{horizontal and vertical parts of the launch velocity}{\si{m.s^{-1}}}
  \end{symbols}
  \whentouse{any launch at an angle. Back again: $u = \sqrt{u_x^2 + u_y^2}$,
    $\theta = \tan^{-1}(u_y/u_x)$.}
  \derive{in the right-angled triangle, $\cos\theta = u_x/u$ and $\sin\theta = u_y/u$.}
  \trap{use degree mode. If the angle is from the \emph{vertical}, $u_x = u\sin\theta$.}
\end{formulacard}

\textbf{The equations of motion} (constant acceleration only):
\[ v = u + at \qquad s = ut + \frac{1}{2}at^2 \qquad v^2 = u^2 + 2as \]
{\centering\hscbadgeonsheet\par}
\begin{itemize}
  \item $u$, $v$: initial and final velocity (\si{m.s^{-1}}); $a$: acceleration (\si{m.s^{-2}});
        $s$: \term{displacement}, the change in position, with a sign (\si{m}); $t$: time (\si{s}).
  \item Use them once for each direction. Horizontally $a_x = 0$, so $v_x = u_x$ and
        $s_x = u_x t$. Vertically $a_y = \qty{-9.8}{m.s^{-2}}$. Never mix $x$ and $y$ values.
\end{itemize}

\textbf{Graphs against time:} $s_x$--$t$ is a straight line through the origin; $v_x$--$t$ is
flat; $s_y$--$t$ is an upside-down parabola; $v_y$--$t$ is a straight line with gradient
\qty{-9.8}{m.s^{-2}}; $a_y$--$t$ is flat at \qty{-9.8}{m.s^{-2}}.

\subsection{The simplest case: a horizontal launch}
Off a table, a ledge or a cliff, $u_y = 0$: the fall is like a drop from rest. So the time in
the air depends \textbf{only on the height}. Four times the height doubles the time (and, at
the same speed, the range).

\begin{formulacard}[notonsheet]{Horizontal launch from a height}
  \cardformula{t = \sqrt{\frac{2h}{g}} \qquad R = u\sqrt{\frac{2h}{g}}}
  \begin{symbols}
    \sym{$t$, $h$}{time of flight; launch height above the landing point}{\si{s}; \si{m}}
    \sym{$u$}{horizontal launch speed}{\si{m.s^{-1}}}
    \sym{$g$}{size of gravity's acceleration, \num{9.8}}{\si{m.s^{-2}}}
    \sym{$R$}{range, from the point under the launch point}{\si{m}}
  \end{symbols}
  \whentouse{the launch velocity is horizontal ($u_y = 0$).}
  \derive{$s = ut + \tfrac{1}{2}at^2$ with $s_y = -h$, $u_y = 0$, $a_y = -g$ gives
    $-h = -\tfrac{1}{2}gt^2$, so $t = \sqrt{2h/g}$. Then $R = u_x t$.}
  \trap{use $u_y = 0$, not $u_y = u$. Here $h$ is positive: do not put in $-h$.}
\end{formulacard}

%% ==================================================================
\section{Height, time of flight, range and final velocity}\label{ch1:sec:height}
\covers{M5.IQ1.02, M5.IQ1.04}

\begin{bigidea}
At the top of the path the vertical velocity is zero for an instant. This gives the maximum
height and, on level ground, the time of flight and the range.
\end{bigidea}

\begin{keywords}
  \kw{Maximum height ($h_\mathrm{max}$)}{the greatest height above the launch point.}
  \kw{Time of flight ($t_\mathrm{flight}$)}{the time from launch until landing.}
  \kw{Range ($R$)}{the horizontal distance from the launch point to the landing point.}
\end{keywords}

\begin{sidebyside}[0.58]
  \centering
  \begin{tikzpicture}[x=0.7cm, y=0.7cm,
      comp/.style={->, line width=0.8pt, color=hscvel!55}]
    \draw[hscink, line width=0.6pt] (-0.4,0) -- (9.4,0);
    \draw[hscmuted!60, line width=1.1pt, domain=0:2.449, samples=80, variable=\t]
      plot ({0.37*9*\t},{0.37*(12*\t-4.9*\t*\t)});
    \draw[comp] (0,0) -- (1.08,0);
    \draw[comp] (0,0) -- (0,1.44) node[left, font=\footnotesize, text=hscvel] {$u_y$};
    \draw[velocity] (0,0) -- (1.08,1.44) node[above, text=hscvel, inner sep=1pt] {$u$};
    \node[font=\footnotesize, text=hscvel, below] at (0.6,0) {$u_x$};
    \draw[comp] (2.039,2.039) -- (3.119,2.039);
    \draw[comp] (2.039,2.039) -- (2.039,2.759);
    \draw[velocity] (2.039,2.039) -- (3.119,2.759);
    \draw[velocity] (4.078,2.718) -- (5.158,2.718)
      node[right, text=hscvel, inner sep=2pt, font=\footnotesize] {$v = v_x$, \ $v_y = 0$};
    \draw[accel] (4.078,2.56) -- (4.078,1.80)
      node[below, text=hscacc, font=\footnotesize, inner sep=1pt] {$a = g$ (down)};
    \draw[comp] (6.117,2.039) -- (7.197,2.039);
    \draw[comp] (6.117,2.039) -- (6.117,1.319);
    \draw[velocity] (6.117,2.039) -- (7.197,1.319);
    \draw[comp] (8.156,0) -- (9.236,0);
    \draw[comp] (8.156,0) -- (8.156,-1.44) node[left, font=\footnotesize, text=hscvel] {$v_y$};
    \draw[velocity] (8.156,0) -- (9.236,-1.44) node[right, text=hscvel, inner sep=1pt] {$v$};
    \foreach \p in {(0,0),(2.039,2.039),(4.078,2.718),(6.117,2.039),(8.156,0)}
      \fill[object] \p circle (3.2pt);
    \draw[guide] (8.4,2.718) -- (10.05,2.718);
    \draw[displacement, <->] (9.9,0) -- (9.9,2.718)
      node[midway, right, text=hscink, font=\small] {$h_\mathrm{max}$};
    \draw[<->, line width=0.6pt, hscink] (0,-1.75) -- (8.156,-1.75)
      node[midway, fill=white, font=\small, inner sep=2pt] {range $R$};
    \node[font=\footnotesize, text=hscmuted] at (2.0,-0.55) {time up};
    \node[font=\footnotesize, text=hscmuted] at (6.1,-0.55) {time down};
  \end{tikzpicture}
  \captionof{figure}{Velocity on level ground. Light arrows: $v_x$ and $v_y$. At the top only
    $v_x$ is left ($v_y = 0$).}\label{ch1:fig:path}
\nextside
  \begin{itemize}
    \item At the top \textbf{only $v_y = 0$}. The ball still moves sideways at $v_x$.
    \item The acceleration is still $g$ down, at every point. It is never zero.
    \item Level ground: the path is \term{symmetrical}; time up $=$ time down.
  \end{itemize}
\end{sidebyside}

\begin{formulacard}[notonsheet]{Maximum height, time of flight and range}
  \cardformula{t_\mathrm{top} = \frac{u_y}{g} \quad\; h_\mathrm{max} = \frac{u_y^2}{2g}
    \quad\; t_\mathrm{flight} = \frac{2u_y}{g} \quad\; R = \frac{u^2\sin 2\theta}{g}}
  \begin{symbols}
    \sym{$t_\mathrm{top}$}{time from launch to the top}{\si{s}}
    \sym{$t_\mathrm{flight}$}{time from launch to landing}{\si{s}}
    \sym{$h_\mathrm{max}$}{greatest height above the launch point}{\si{m}}
    \sym{$R$}{range on level ground}{\si{m}}
    \sym{$u_y$, $u$}{vertical part of the launch velocity; launch speed}{\si{m.s^{-1}}}
    \sym{$\theta$}{launch angle above the horizontal}{\si{\degree}}
    \sym{$g$}{size of gravity's acceleration, \num{9.8}}{\si{m.s^{-2}}}
  \end{symbols}
  \whentouse{$t_\mathrm{top}$ and $h_\mathrm{max}$: any upward launch. $t_\mathrm{flight}$ and
    $R$: launch and landing at the same height, no air resistance.}
  \derive{put $a_y = -g$ into the sheet equations. At the top $v_y = 0$: $v = u + at$ gives
    $0 = u_y - g\,t_\mathrm{top}$, and $v^2 = u^2 + 2as$ gives $0 = u_y^2 - 2g\,h_\mathrm{max}$.
    At landing $s_y = 0$: $s = ut + \tfrac{1}{2}at^2$ gives $0 = t\left(u_y - \tfrac{1}{2}gt\right)$,
    so $t_\mathrm{flight} = 2u_y/g$ ($t = 0$ is the launch). Then
    $R = u_x t_\mathrm{flight} = u\cos\theta \times \dfrac{2u\sin\theta}{g} = \dfrac{u^2\sin 2\theta}{g}$,
    because $2\sin\theta\cos\theta = \sin 2\theta$.}
  \trap{use $u_y$, not $u$. Put in $g = 9.8$, not $-9.8$ (the minus sign is built in).}
\end{formulacard}

\begin{itemize}
  \item With $u_y = u\sin\theta$: \ $h_\mathrm{max} = \dfrac{u^2\sin^2\theta}{2g}$ and
        $t_\mathrm{flight} = \dfrac{2u\sin\theta}{g}$. Double $u$ (same angle): the times double;
        the height and the range become 4 times bigger.
\end{itemize}

\subsection{Range and the \texorpdfstring{\ang{45}}{45 degree} rule}
\begin{itemize}
  \item The range needs horizontal speed \emph{and} time in the air. A low launch lands
        quickly; a steep one moves forwards slowly.
  \item \textbf{\ang{45} gives the longest range}, because $\sin 2\theta = 1$ when
        $2\theta = \ang{90}$. Then $R_\mathrm{max} = u^2/g$. This is true only on level ground
        with no air resistance.
  \item \textbf{\term{Complementary angles}} (two angles that add to \ang{90}), such as \ang{30}
        and \ang{60}, give the same range. The steeper one goes higher and stays up longer.
\end{itemize}

\subsection{Launching from a height}
\begin{itemize}
  \item \term{Launch height} $h$: launch point above landing point, so $s_y = -h$ at landing.
  \item Solve $-h = u_y t + \frac{1}{2}a_y t^2$ for $t$. It is a \term{quadratic equation}
        (an equation with a $t^2$ term, a $t$ term and a number). Write it as
        $4.9t^2 - u_y t - h = 0$, so $A = 4.9$, $B = -u_y$, $C = -h$ in
        $t = \dfrac{-B \pm \sqrt{B^2 - 4AC}}{2A}$. A \term{root} is a solution. Keep the
        \textbf{positive} root; the negative one is before the launch. Then $R = u_x t$.
  \item The level-ground $t_\mathrm{flight}$ and $R$ do not apply. A target \emph{above} the
        launch point gives two positive roots (way up, way down).
  \item Higher launch point: longer time and range, same $h_\mathrm{max}$ above the launch point.
  \item From a height the best angle is \textbf{less than \ang{45}}. The drop adds time in the
        air, so it is better to put more of the speed into $u_x$.
\end{itemize}

\subsection{Final velocity}
The \term{final velocity} is the velocity just before landing: a speed \emph{and} a direction.
\begin{itemize}
  \item $v_x = u_x$. Then $v_y = u_y + a_y t$ or $v_y^2 = u_y^2 + 2a_y s_y$ ($v_y$ is negative
        when the ball moves down).
  \item Speed $v = \sqrt{v_x^2 + v_y^2}$; direction $\tan\alpha = \lvert v_y\rvert/v_x$, where
        $\alpha$ is the angle below the horizontal (not on the formula sheet).
  \item On \textbf{level ground}, $s_y = 0$ gives $v_y = -u_y$. The final speed equals the launch
        speed, at the launch angle but below the horizontal.
\end{itemize}

%% ==================================================================
\section{Energy and air resistance}\label{ch1:sec:energy}
\covers{M5.IQ1.01, M5.IQ1.02}

\begin{bigidea}
With no air resistance, kinetic plus potential energy stays the same during the flight. This
gives the speed at any height, without the time or the angle.
\end{bigidea}

\begin{itemize}
  \item On the way up, \term{kinetic energy} ($K = \frac{1}{2}mv^2$) turns into
        \term{gravitational potential energy} ($\Delta U = mg\Delta h$). On the way down it
        turns back. The total stays the same: \term{conservation of energy}.
  \item At the top of an angled flight, $K$ is smallest but \emph{not} zero:
        $K_\mathrm{top} = \frac{1}{2}mu_x^2$.
\end{itemize}

Kinetic energy gained $=$ potential energy lost, so (the mass cancels):
\[ \frac{1}{2}mv_2^2 - \frac{1}{2}mv_1^2 = mg(h_1 - h_2) \quad\Longrightarrow\quad
   v_2^2 = v_1^2 + 2g(h_1 - h_2) \]
{\centering\hscbadgecombine{$K = \frac{1}{2}mv^2$ and $\Delta U = mg\Delta h$}\par}
\begin{itemize}
  \item $v_1$, $v_2$ are the speeds at heights $h_1$, $h_2$ (\si{m.s^{-1}}, \si{m}). Landing $h$
        below the launch: $v = \sqrt{u^2 + 2gh}$. $h_1 - h_2$ is negative when the ball ends
        higher.
  \item Energy has no direction: it gives only the speed. For the angle, use $v_x = u_x$ and
        $\cos\alpha = v_x/v$.
\end{itemize}

\subsection{Air resistance: where the model fails}
\term{Drag} (air resistance) is an extra force opposite to the velocity. It grows quickly with
speed. Compared with the model:
\begin{sidebyside}[0.5]
  \begin{itemize}
    \item the \textbf{range is shorter} (drag slows the horizontal motion);
    \item the \textbf{maximum height is lower} (on the way up, drag adds to gravity);
  \end{itemize}
\nextside
  \begin{itemize}
    \item the path is \textbf{not symmetrical}: the way down is steeper and takes longer;
    \item the ball \textbf{lands slower}, and the best angle is \textbf{below \ang{45}}.
  \end{itemize}
\end{sidebyside}

The model works \textbf{well} for small, \term{dense} (heavy for their size), slow objects over
short distances: drag hardly changes a heavy object's motion. It works \textbf{badly} for light,
large or fast objects (a feather, a golf ball after a drive).

%% ==================================================================
\section{Predicting motion: uniform velocity and constant acceleration}\label{ch1:sec:predict}
\covers{M5.IQ1.04}

Uniform velocity and constant acceleration are the two models of straight-line motion. A
projectile uses both at once, with the same time $t$. A \term{prediction} is a result
calculated from a model before it is measured.

\begin{itemize}
  \item \term{Uniform velocity} (the same speed in the same direction): zero net force
        (Newton's first law). $s = ut$ and $v = u$ (the sheet equations with $a = 0$). The
        displacement--time graph is a straight line; its gradient is the velocity.
  \item \term{Constant acceleration} (the velocity changes by the same amount every second): a
        constant net force ($\vec{F}_\mathrm{net} = m\vec{a}$), as in \term{free fall} (gravity
        only). The velocity--time graph is a straight line: gradient $=$ acceleration, area
        under it $=$ displacement.
  \item \textbf{Choosing the model}, for each direction: no net force, uniform velocity; a constant net
        force, constant acceleration; a changing force (large drag), neither is exact. Never
        use $s = ut$ for a falling object.
  \item \textbf{Straight up} ($u_x = 0$): time up $=$ time down; the ball returns at its launch
        speed. Below the top it passes each height twice (two positive times).
  \item \textbf{From a moving vehicle:} a released object keeps the vehicle's velocity
        (\term{inertia}). A ball tossed straight up from a cart at uniform velocity lands back
        in the cart; seen from the ground, its path is a parabola. A parcel dropped from a plane
        at uniform velocity stays below the plane.
  \item \textbf{To predict:} one origin, a model for each object and direction, and one clock
        $t$. ``At the same place'' gives one equation in $t$.
\end{itemize}

%% ==================================================================
\section{Key experiments}\label{ch1:sec:experiments}
\covers{M5.IQ1.03}

\begin{experiment}{Range of a ball launched horizontally}
\aim{To collect \term{primary data} (data you measure) to \term{validate} (test) $R = u\sqrt{2h/g}$:
  $R$ against $u$ should be a straight line through the origin, with gradient $\sqrt{2h/g}$.}
\method{\begin{sidebyside}[0.5]
  \centering
  \begin{tikzpicture}[x=0.62cm, y=0.62cm]
      \fill[ground] (-4.6,-0.22) rectangle (4.8,0);
      \draw[hscink, line width=0.6pt] (-4.6,0) -- (4.8,0);
      \fill[hscmuted!22] (-4.4,0) rectangle (-0.12,3.0);
      \draw[hscmuted, line width=0.6pt] (-4.4,0) rectangle (-0.12,3.0);
      \fill[hscmuted!40] (-3.8,3.0) -- (-3.8,5.0) .. controls (-2.8,3.2) and (-1.8,3.18)
        .. (0,3.18) -- (0,3.0) -- cycle;
      \draw[hscink, line width=1.2pt] (-3.8,5.0) .. controls (-2.8,3.2) and (-1.8,3.18) .. (0,3.18);
      \fill[object] (-3.22,4.42) circle (3pt);
      \fill[object] (0,3.33) circle (3pt);
      \draw[hscink, line width=1pt] (-0.62,3.18) -- (-0.62,3.85) -- (-0.22,3.85) -- (-0.22,3.18);
      \draw[hscforce, line width=0.6pt, dashed] (-0.62,3.5) -- (-0.22,3.5);
      \draw[guide, line width=0.8pt] (0,3.33) .. controls (1.3,3.33) and (2.4,2.0) .. (2.9,0.15);
      \fill[object] (2.9,0.15) circle (3pt);
      \draw[hscink, line width=0.4pt] (0.06,3.12) -- (0.06,0.3);
      \fill[hscink] (0.06,0.3) -- (-0.02,0.42) -- (0.14,0.42) -- cycle;
      \fill[hscink!85] (2.3,0) rectangle (3.6,0.07);
      \draw[guide] (0.2,3.18) -- (4.5,3.18);
      \draw[displacement, <->] (4.3,0) -- (4.3,3.18) node[midway, right, text=hscink, font=\small] {$h$};
      \draw[<->, line width=0.6pt, hscink] (0.06,-0.6) -- (2.9,-0.6)
        node[midway, below, font=\small] {$R$};
      \node[diagramlabel, font=\footnotesize] at (-2.3,1.5) {bench};
      \node[diagramlabel, font=\footnotesize, anchor=south] (lg) at (0.6,4.6) {light gate};
      \draw[hscmuted, line width=0.4pt] (lg.south) -- (-0.42,3.88);
      \node[diagramlabel, font=\footnotesize, anchor=west, align=left] (pl) at (0.45,1.2) {plumb\\line};
      \draw[hscmuted, line width=0.4pt] (pl.west) -- (0.1,1.2);
      \node[diagramlabel, font=\footnotesize, anchor=north west, align=left] (cp) at (3.25,-0.25) {carbon\\paper};
      \draw[hscmuted, line width=0.4pt] (cp.north west) -- (3.3,0.05);
  \end{tikzpicture}
  \captionof{figure}{Side view of the bench set-up.}\label{ch1:fig:apparatus}
\nextside
  \begin{itemize}
    \item A steel ball leaves the ramp's horizontal end, so $u_y = 0$.
    \item A \term{light gate} times the ball blocking its beam: $u = \text{diameter} \div \Delta t$.
    \item Carbon paper marks the landing; $R$ is measured from the plumb line.
    \item Change $u$ (5 speeds, 5 trials each); keep $h$ and the ball the same.
  \end{itemize}
\end{sidebyside}}
\results{$R$ against $u$ ($h = \qty{0.900}{m}$) is a straight line very close to the origin, with
  gradient \qty{0.423}{s} (model: \qty{0.429}{s}). The value of $g$ found from the gradient is
  about \qty{3}{\percent} above \qty{9.8}{m.s^{-2}}.}
\conclusion{$R$ is \term{proportional} to $u$ (double $u$ and $R$ doubles): the data validate
  the model for a dense, slow
  ball. Every $R$ is about \qty{2}{\percent} low: a \term{systematic error} (same direction each
  time), not a model failure.}
\significance{The fall time is the same at every speed, so $R$ grows in step with $u$: the
  horizontal and vertical motions are independent.}
\end{experiment}

% kept together (a minipage), so the three bullets do not split across a page break
\par\noindent\begin{minipage}{\linewidth}
\textbf{Other ways to validate the model:}
\begin{itemize}
  \item $R$ against $\sin 2\theta$ (fixed $u$, level ground): a straight line through the
        origin. The largest $R$ is at \ang{45}; \ang{30} and \ang{60} give equal $R$.
  \item $R$ against $\sqrt{h}$ (fixed $u$, horizontal launch): a straight line through the origin.
  \item Video analysis: $v_y$--$t$ is a straight line, gradient \qty{-9.8}{m.s^{-2}}; $v_x$--$t$ is flat.
\end{itemize}
\end{minipage}

%% ==================================================================
\begin{summarybox}{Remember}
\begin{itemize}
  \item $u_x = u\cos\theta$, $u_y = u\sin\theta$; $v_x$ is constant; $a_y = \qty{-9.8}{m.s^{-2}}$ even
        at the top.
  \item Horizontal launch: $t = \sqrt{2h/g}$. The fall time depends only on the height.
  \item Level ground: $h_\mathrm{max} = u_y^2/2g$, $t_\mathrm{flight} = 2u_y/g$,
        $R = u^2\sin 2\theta/g$ (longest at \ang{45}).
  \item From a height: solve $-h = u_y t - \tfrac{1}{2}gt^2$; keep the positive root.
  \item A final velocity needs a size \emph{and} a direction; energy gives only the speed.
  \item Drag: shorter, lower, steeper path. The model suits dense, slow objects.
\end{itemize}
\end{summarybox}
}{}% Projectile Motion
\IfFileExists{./ch2.tex}{\setcounter{chapter}{1}% ch2.tex -- Module 5 ESSENTIALS, Chapter 2: Circular Motion: Forces and Real Situations
% Condensed from the final full book: src/mod5/ch2.tex (same notation, same facts).
% Knowledge points only: no worked examples, no questions, no answers.
% All diagrams are the full book's TikZ code, copied (no figure files are needed).
\chapter[Circular Motion: Forces]{Circular Motion: Forces and Real Situations}\label{ch2:chap}
\covers{M5.IQ2.01, M5.IQ2.02, M5.IQ2.03}

\begin{summarybox}[unbreakable]{In this chapter}
\begin{itemize}
  \item Why moving in a circle at a constant speed is accelerating.
  \item The formulas for $v$, $a_\mathrm{c}$, $F_\mathrm{c}$ and $\omega$.
  \item Which real force provides the centripetal force in each situation.
  \item The whirling-stopper investigation of $F_\mathrm{c}$, $m$, $v$ and $r$.
\end{itemize}
\end{summarybox}

%% ==================================================================
\section{Uniform circular motion and centripetal force}\label{ch2:sec:ucm}
\covers{M5.IQ2.01, M5.IQ2.02, M5.IQ2.03}

\begin{bigidea}
An object moving in a circle at a constant speed is accelerating, because its direction
keeps changing. The acceleration points to the centre. A net force towards the centre,
the centripetal force, causes it.
\end{bigidea}

\begin{keywords}
  \kw{Uniform circular motion}{moving around a circle at a constant speed.}
  \kw{Period ($T$)}{the time for one full turn. Unit: second (s).}
  \kw{Tangent}{a line that touches the circle at one point, at \ang{90} to the radius.}
  \kw{Centripetal acceleration ($a_\mathrm{c}$)}{the acceleration towards the centre of the circle.}
  \kw{Centripetal force ($F_\mathrm{c}$)}{the net force towards the centre of the circle.}
\end{keywords}

\begin{itemize}
  \item Velocity is a \term{vector}: it has a size (the speed) \emph{and} a direction. In a
        circle the direction changes all the time. So the velocity changes: the object
        accelerates.
  \item The velocity is always along the tangent (\cref{ch2:fig:ucm}a). Two velocities a
        moment apart, drawn tail to tail, give a change $\Delta v$ that points to the centre
        (\cref{ch2:fig:ucm}b).
  \item So the acceleration $a_\mathrm{c}$ points to the centre (``centripetal'' means
        ``centre-seeking''). One lap is the \term{circumference} $2\pi r$, done in one
        period $T$, so $v = 2\pi r/T$.
  \item $a_\mathrm{c}$ grows with $v^2$: a faster object turns its velocity arrow more
        quickly, \emph{and} the arrow is longer. At the same speed, a smaller $r$ gives a bigger
        $a_\mathrm{c}$.
  \item By Newton's second law ($\vec{F}_\mathrm{net} = m\vec{a}$), the \term{net force} (the
        total of all the forces) also points to the centre: $F_\mathrm{net} = ma_\mathrm{c} = mv^2/r$.
\end{itemize}

\begin{figure}
  \centering
  \begin{tikzpicture}[x=1cm, y=1cm]
    \begin{scope}
      \draw[hscmuted!55, line width=0.8pt] (0,0) circle (1.4);
      \fill[hscink] (0,0) circle (1.5pt);
      \node[diagramlabel, below, inner sep=3pt] at (0,0) {centre};
      \draw[->, hscmuted, line width=0.6pt] (212:1.75) arc (212:255:1.75);
      \node[diagramlabel, text=hscmuted, anchor=north east] at (232:1.85) {motion};
      \draw[guide] (35:1.4) -- ++(125:1.9);
      \node[diagramlabel, text=hscmuted, anchor=south west, inner sep=1pt] at ($(35:1.4)+(125:1.9)+(-0.1,0.02)$)
        {path if the string breaks};
      \foreach \a in {35, 155, 275} {
        \coordinate (P) at (\a:1.4);
        \draw[velocity] (P) -- ++(\a+90:1.0);
        \draw[accel] (P) -- ($(P)!0.5!(0,0)$);
        \fill[object] (P) circle (4pt);
      }
      \node[diagramlabel, text=hscvel, anchor=west] at ($(35:1.4)+(125:0.9)+(0.08,0)$) {$v$};
      \node[diagramlabel, text=hscacc, anchor=north west, inner sep=1pt] at ($(35:0.85)+(0.05,-0.02)$) {$a_\mathrm{c}$};
      \node[diagramlabel, text=hscvel] at ($(155:1.4)+(245:1.2)$) {$v$};
      \node[diagramlabel, text=hscacc, anchor=south, inner sep=2pt] at ($(155:0.85)+(0.05,0.05)$) {$a_\mathrm{c}$};
      \node[diagramlabel, text=hscvel, anchor=west] at ($(275:1.4)+(5:1.05)$) {$v$};
      \node[diagramlabel, text=hscacc, anchor=west, inner sep=2pt] at ($(275:0.95)+(0.02,-0.1)$) {$a_\mathrm{c}$};
      \node[font=\small\bfseries, text=hscink] at (-1.8,1.8) {(a)};
    \end{scope}
    \begin{scope}[xshift=6.4cm, yshift=-0.8cm]
      \coordinate (O) at (0,0);
      \coordinate (A) at (110:1.8);
      \coordinate (B) at (160:1.8);
      \draw[velocity] (O) -- (A) node[pos=0.55, right, inner sep=2pt, text=hscvel] {$v_1$ (earlier)};
      \draw[velocity] (O) -- (B) node[pos=0.45, below left, inner sep=2pt, text=hscvel] {$v_2$ (later)};
      \draw[accel] (A) -- (B) node[pos=0.5, above left, inner sep=1pt, text=hscacc] {$\Delta v$};
      \fill[hscink] (O) circle (1.3pt);
      \node[font=\small\bfseries, text=hscink] at (-2.1,2.5) {(b)};
    \end{scope}
  \end{tikzpicture}
  \caption{(a) Top view of a ball whirled on a string: $v$ is along the tangent,
    $a_\mathrm{c}$ points to the centre. (b) Two velocities a moment apart, tail to tail:
    $\Delta v$ points to the centre.}
  \label{ch2:fig:ucm}
\end{figure}

\begin{formulacard}{Speed, acceleration and force in a circle}
  \cardformula{v = \frac{2\pi r}{T} \qquad\quad a_\mathrm{c} = \frac{v^2}{r} \qquad\quad
    F_\mathrm{c} = \frac{mv^2}{r}}
  \begin{symbols}
    \sym{$v$}{speed of the object (along the tangent)}{\si{m.s^{-1}}}
    \sym{$r$}{radius of the circle}{\si{m}}
    \sym{$T$}{period: the time for one full turn}{\si{s}}
    \sym{$a_\mathrm{c}$}{centripetal acceleration, towards the centre}{\si{m.s^{-2}}}
    \sym{$F_\mathrm{c}$}{centripetal force: the net force towards the centre}{\si{N}}
    \sym{$m$}{mass of the object}{\si{kg}}
  \end{symbols}
  \whentouse{any object moving in a circle at a constant speed.}
  \trap{$T$ is the period (not tension); use the radius, not the diameter. $F_\mathrm{c}$ is
    the net force: never draw it as an extra arrow.}
\end{formulacard}

\begin{itemize}
  \item $F_\mathrm{c} \propto m$, \ $F_\mathrm{c} \propto v^2$, \ $F_\mathrm{c} \propto 1/r$
        ($\propto$ means ``is proportional to''): double $v$ and $F_\mathrm{c}$ is 4 times
        bigger. With the period, $a_\mathrm{c} = 4\pi^2 r/T^2$ (combine $v = 2\pi r/T$ and
        $a_\mathrm{c} = v^2/r$).
  \item \term{Frequency} $f$ (turns each second, in \si{Hz}): $f = 1/T$ (on the sheet).
        Convert first: \si{km.h^{-1}}~$\div\, 3.6$ gives \si{m.s^{-1}}; revolutions per
        minute (rpm)~$\div\, 60$ gives $f$.
  \item \term{Angular velocity} $\omega$ is the angle turned each second:
        $\omega = \Delta\theta/t$ (on the sheet), in \si{rad.s^{-1}}. One full turn is $2\pi$
        \term{radians} (rad) $= \ang{360}$, so $\omega = 2\pi/T$ (not on the sheet).
  \item The centripetal force is \emph{not} a new kind of force: real forces do this job.
        The real force can be tension (ball on a string), friction (flat bend), the normal
        force (banked track, rotor ride) or gravity (Moon or satellite in orbit).
  \item No force pushes you outwards on a turn. By \term{inertia} (the tendency to keep
        moving in a straight line) your body tries to go straight; the seat, seatbelt and
        door push you \emph{inwards}. ``Centrifugal force'' is not a real force.
  \item If a whirling ball's string breaks, the inward force is gone. By Newton's first law,
        the ball moves in a straight line along the tangent, not straight outwards.
\end{itemize}

%% ==================================================================
\section{Real situations: the forces that provide \texorpdfstring{$F_\mathrm{c}$}{Fc}}\label{ch2:sec:situations}
\covers{M5.IQ2.02, M5.IQ2.03}

\begin{bigidea}
In every situation, find the real forces and split them into a part towards the centre
and a vertical part. The part towards the centre adds up to $mv^2/r$.
\end{bigidea}

\begin{keywords}
  \kw{Tension ($F_\mathrm{T}$)}{the pull of a string, rope or chain, along its length.}
  \kw{Normal force ($F_\mathrm{N}$)}{the push of a surface on an object, at \ang{90} to the surface.}
  \kw{Component}{the part of a force that acts in one chosen direction.}
\end{keywords}

\textbf{Method.} Draw a \term{free-body diagram}: the object with one arrow per real force (no
$F_\mathrm{c}$ arrow).
With $\theta$ from the \emph{vertical}, a force has parts $F\cos\theta$ (vertical) and
$F\sin\theta$ (horizontal). Towards the centre they add up to $mv^2/r$; vertically
(horizontal circle) to zero.

\subsection*{Strings: conical pendulums and swings}
A \term{conical pendulum} is a mass on a string moving in a horizontal circle. Only
$F_\mathrm{T}$ and $mg$ act on it (\cref{ch2:fig:conical}).

\begin{figure}
  \centering
  \begin{tikzpicture}[x=1cm, y=1cm]
    \begin{scope}
      \fill[ground] (-1.5,0) rectangle (2.4,0.22);
      \draw[hscink, line width=0.8pt] (-1.5,0) -- (2.4,0);
      \coordinate (piv) at (0,0);
      \coordinate (B) at (1.548,-2.212);
      \coordinate (C) at (0,-2.212);
      \draw[guide] (piv) -- (0,-2.7);
      \draw[hscmuted, dashed, line width=0.7pt] (C) ellipse (1.548 and 0.36);
      \draw[hscink, line width=1pt] (piv) -- (B) node[pos=0.5, right, diagramlabel] {$L$};
      \draw[hscink, line width=0.5pt] (0,-0.7) arc (270:305:0.7);
      \node[diagramlabel] at (0.19,-0.92) {$\theta$};
      \draw[displacement, <->] (C) -- (B) node[pos=0.5, below, diagramlabel, inner sep=2pt] {$r$};
      \fill[hscink] (C) circle (1.5pt);
      \draw[hscmuted, line width=0.4pt] (C) -- (-1.75,-2.75);
      \node[diagramlabel, anchor=north east, align=right, inner sep=1pt] at (-1.75,-2.75) {centre of\\the circle};
      \fill[object] (B) circle (5pt);
      \fill[hscink] (piv) circle (1.5pt);
      \node[diagramlabel, anchor=east] at (-0.12,-0.25) {pivot};
      \node[font=\small\bfseries, text=hscink] at (-1.3,-0.75) {(a)};
    \end{scope}
    \begin{scope}[xshift=7.4cm, yshift=-1.7cm]
      \coordinate (O) at (0,0);
      \draw[guide] (O) -- (0,1.5);
      \draw[guide] (0,1.5) -- (-1.05,1.5);
      \draw[force] (O) -- (-1.05,1.5) node[pos=0.62, below left, inner sep=1pt, text=hscforce] {$F_\mathrm{T}$};
      \draw[force] (O) -- (0,-1.5) node[pos=0.8, right, text=hscforce] {$mg$};
      \draw[hscink, line width=0.5pt] (0,0.55) arc (90:125:0.55);
      \node[diagramlabel] at (-0.15,0.8) {$\theta$};
      \node[diagramlabel, anchor=west] at (0.05,1.05) {$F_\mathrm{T}\cos\theta$};
      \node[diagramlabel, anchor=south] at (-0.53,1.53) {$F_\mathrm{T}\sin\theta$};
      \draw[displacement] (-1.4,-0.6) -- (-2.5,-0.6);
      \node[diagramlabel, anchor=north, text=hscmuted] at (-1.95,-0.65) {to the centre};
      \fill[object] (O) circle (5pt);
      \node[font=\small\bfseries, text=hscink] at (-2.5,1.8) {(b)};
    \end{scope}
  \end{tikzpicture}
  \caption{A conical pendulum. (a) Side view: the centre of the circle is directly below
    the pivot. (b) Free-body diagram; dashed lines are the components of the tension.}
  \label{ch2:fig:conical}
\end{figure}

\begin{itemize}
  \item The centre of the circle is directly \emph{below} the \term{pivot} (the fixed point),
        so $r = L\sin\theta$. Towards the centre: $F_\mathrm{T}\sin\theta = mv^2/r$.
        Vertically: $F_\mathrm{T}\cos\theta = mg$, so $F_\mathrm{T} = mg/\cos\theta > mg$.
  \item Swing-ride chair (hangs from a spinning frame): $r = d + L\sin\theta$, $d$ = axis to chain top.
  \item A \term{vertical circle} (an upright circle) is \emph{not} uniform: the speed changes.
        At each point the net force to the centre is still $mv^2/r$. Lowest point of a swing:
        $F_\mathrm{T} - mg = mv^2/r$ (you feel heavier). Top, on a string:
        $F_\mathrm{T} + mg = mv^2/r$. Neither is on the sheet.
\end{itemize}

\subsection*{Tilted surfaces: bowls and banked tracks}
\begin{itemize}
  \item A \term{banked track} is higher on the outside of the bend. Its \term{bank angle}
        $\theta$ is its tilt from the horizontal. $F_\mathrm{N}$ (at \ang{90} to the road)
        tilts to the centre, by the same $\theta$ from the vertical.
  \item At the \term{design speed} no friction is needed. Faster: the car tends to slide up
        the slope, so friction acts \emph{down} the slope. Slower: friction acts \emph{up}
        the slope (\cref{ch2:fig:bankfriction}).
\end{itemize}

\begin{figure}
  \centering
  % full book figure ch2:fig:bankfriction (bank angle 20 deg). Added here: the angle theta
  % between F_N and the vertical in (b) (F_N is drawn at 110 deg, so 20 deg from the vertical).
  \begin{tikzpicture}[x=1.2cm, y=1.2cm]
    \foreach \dx in {0, 4.0, 8.0} {
      \begin{scope}[xshift=\dx*1.2cm]
        \fill[ground] (0,0) -- (2.8,1.019) -- (2.8,0.70) -- (0,-0.319) -- cycle;
        \draw[hscink, line width=1pt] (0,0) -- (2.8,1.019);
        \draw[guide] (0,0) -- (0.9,0);
        \draw[hscink, line width=0.5pt] (0.5,0) arc (0:20:0.5);
        \node[diagramlabel, font=\footnotesize] at (0.66,0.1) {$\theta$};
        \begin{scope}[rotate around={20:(1.4,0.510)}]
          \draw[object, rounded corners=1.5pt] (0.98,0.60) rectangle (1.82,0.98);
          \fill[hscink] (1.15,0.60) circle (0.085);
          \fill[hscink] (1.65,0.60) circle (0.085);
        \end{scope}
        \draw[force] (1.314,0.744) -- ++(0,-1.0) node[pos=0.9, right, text=hscforce, inner sep=2pt] {$mg$};
        \draw[force] (1.314,0.744) -- ++(110:1.064) node[pos=1, left, text=hscforce] {$F_\mathrm{N}$};
        \fill[hscink] (1.314,0.744) circle (1.3pt);
      \end{scope}
    }
    % (b): angle of F_N from the vertical
    \draw[guide] (5.314,0.744) -- ++(0,1.1);
    \draw[hscink, line width=0.5pt] (5.314,1.304) arc (90:110:0.56);
    \node[diagramlabel, font=\footnotesize] at ($(5.314,0.744)+(100:0.82)$) {$\theta$};
    % friction: (a) up the slope, (c) down the slope
    \draw[force] (1.314,0.744) -- ++(20:0.75) node[pos=1, above, text=hscforce] {$F_\mathrm{f}$};
    \draw[force] (9.314,0.744) -- ++(200:0.75) node[pos=1, above left, text=hscforce, inner sep=1pt] {$F_\mathrm{f}$};
    \node[diagramlabel, font=\footnotesize, align=center, anchor=north] at (1.4,-0.45)
      {\textbf{(a) slower:} friction up the slope};
    \node[diagramlabel, font=\footnotesize, align=center, anchor=north] at (5.4,-0.45)
      {\textbf{(b) design speed:} no friction};
    \node[diagramlabel, font=\footnotesize, align=center, anchor=north] at (9.4,-0.45)
      {\textbf{(c) faster:} friction down the slope};
  \end{tikzpicture}
  \caption{A car on a banked track, from behind (centre of the bend to the left). Not to
    scale.}\label{ch2:fig:bankfriction}
\end{figure}

\begin{itemize}
  \item At the design speed: $F_\mathrm{N}\sin\theta = mv^2/r$ and $F_\mathrm{N}\cos\theta = mg$,
        so $F_\mathrm{N} = mg/\cos\theta > mg$. Banking means less friction is needed, so wet
        turns are safer.
  \item A marble circling inside a round bowl: the same, with $F_\mathrm{N}$ pointing to the
        bowl's centre. $r$ is the horizontal distance to the centre line, not the bowl's radius.
\end{itemize}

\begin{formulacard}[notonsheet]{Tilted string or surface}
  \cardformula{\tan\theta = \frac{v^2}{rg}}
  \begin{symbols}
    \sym{$\theta$}{angle of $F_\mathrm{T}$ or $F_\mathrm{N}$ from the vertical (bank angle)}{\si{\degree}}
    \sym{$v$}{speed (for a banked road: the design speed)}{\si{m.s^{-1}}}
    \sym{$r$}{radius of the horizontal circle}{\si{m}}
    \sym{$g$}{gravitational acceleration, \qty{9.8}{m.s^{-2}}}{\si{m.s^{-2}}}
  \end{symbols}
  \whentouse{conical pendulum; banked road at its design speed; frictionless bowl.}
  \derive{use the two component equations ($F$ is $F_\mathrm{T}$ or $F_\mathrm{N}$).
    \begin{align*}
      F\sin\theta &= \frac{mv^2}{r} && \text{horizontal: net force towards the centre}\\
      F\cos\theta &= mg && \text{vertical: forces balance}\\
      \tan\theta &= \frac{v^2}{rg} && \text{divide the first line by the second; $F$ and $m$ cancel}
    \end{align*}}
  \trap{for a string, $r$ is not the string length: $r = L\sin\theta$. Do not split forces
    along the slope: the acceleration is horizontal.}
\end{formulacard}

\subsection*{Friction: flat bends and the rotor ride}
\begin{itemize}
  \item \term{Static friction} $F_\mathrm{f}$ stops sliding. It is only as big as needed, up to
        $F_\mathrm{f} \le \mu F_\mathrm{N}$ (not on the sheet). The \term{coefficient of
        friction} $\mu$ (no unit) says how strongly the surfaces grip (about 0.8 on a dry road,
        0.1 on ice).
\end{itemize}

\begin{sidebyside}[0.50]
  \centering
  \begin{tikzpicture}[x=1cm, y=1cm]
    \begin{scope}
      \fill[hscmuted!12] (15:1.7) arc (15:115:1.7) -- (115:2.4) arc (115:15:2.4) -- cycle;
      \draw[hscmuted, line width=0.7pt] (15:1.7) arc (15:115:1.7);
      \draw[hscmuted, line width=0.7pt] (15:2.4) arc (15:115:2.4);
      \fill[hscink] (0,0) circle (1.5pt);
      \node[diagramlabel, anchor=north] at (0,-0.05) {centre};
      \coordinate (Q) at (62:2.05);
      \draw[guide] (0,0) -- (Q);
      \draw[velocity] (Q) -- ++(152:1.0) node[above, text=hscvel] {$v$};
      \draw[force] (Q) -- ($(Q)!0.45!(0,0)$) node[pos=1, below right, inner sep=1pt, text=hscforce] {$F_\mathrm{f}$};
      \draw[object, rotate around={152:(Q)}] ($(Q)+(-0.27,-0.15)$) rectangle ($(Q)+(0.27,0.15)$);
      \node[font=\small\bfseries, text=hscink] at (-1.5,2.3) {(a)};
    \end{scope}
    \begin{scope}[xshift=4.0cm, yshift=1.0cm]
      \coordinate (O) at (0,0);
      \draw[force] (O) -- (0,1.1) node[above, text=hscforce] {$F_\mathrm{N}$};
      \draw[force] (O) -- (0,-1.1) node[below, text=hscforce] {$mg$};
      \draw[force] (O) -- (-1.0,0) node[left, text=hscforce] {$F_\mathrm{f}$};
      \fill[object] (O) circle (5pt);
      \node[font=\small\bfseries, text=hscink] at (-0.9,1.3) {(b)};
    \end{scope}
  \end{tikzpicture}
  \captionof{figure}{A car on a flat bend. (a) Top view. (b) Free-body diagram, seen from
    behind (centre to the left).}\label{ch2:fig:flatbend}
\nextside
  \begin{itemize}
    \item \textbf{Flat bend.} The tyres push the road outwards. By Newton's third law, the
          road pushes the tyres \emph{inwards} (\cref{ch2:fig:flatbend}).
    \item This friction is the centripetal force: $F_\mathrm{f} = mv^2/r$. Vertically,
          $F_\mathrm{N} = mg$.
    \item The road can give at most $\mu mg$. If the bend requires more, the car skids
          towards the outside.
  \end{itemize}
\end{sidebyside}

\begin{sidebyside}[0.31]
  \centering
  \begin{tikzpicture}[x=1cm, y=1cm]
    \draw[hscink, line width=1.2pt] (-1.4,-0.2) -- (-1.4,2.6);
    \draw[hscink, line width=1.2pt] (1.4,-0.2) -- (1.4,2.6);
    \draw[hscmuted, line width=0.7pt] (0,2.6) ellipse (1.4 and 0.25);
    \draw[guide] (0,-0.7) -- (0,3.15);
    \node[diagramlabel, anchor=south] at (0,3.15) {axis};
    \draw[->, hscmuted, line width=0.7pt] (-0.5,3.0) arc (200:340:0.55 and 0.18);
    \fill[ground] (-1.3,-0.75) rectangle (1.3,-0.6);
    \draw[hscink, line width=0.7pt] (-1.3,-0.6) -- (1.3,-0.6);
    \node[diagramlabel, anchor=north, text=hscmuted] at (0,-0.8) {floor has dropped};
    \fill[hscbrand!45] (1.15,0.4) rectangle (1.36,1.6);
    \fill[hscbrand!45] (1.255,1.77) circle (0.15);
    \node[diagramlabel, anchor=west] at (1.45,1.75) {rider};
    \node[diagramlabel, anchor=west] at (1.45,2.4) {wall};
    \draw[displacement, <->] (0,0.6) -- (-1.4,0.6) node[pos=0.5, above, diagramlabel] {$r$};
    \coordinate (R) at (1.255,1.05);
    \draw[force] (R) -- (0.1,1.05) node[pos=0.55, above, text=hscforce] {$F_\mathrm{N}$};
    \draw[force] (R) -- (1.255,2.25) node[pos=0.95, left, text=hscforce] {$F_\mathrm{f}$};
    \draw[force] (R) -- (1.255,-0.15) node[pos=0.9, left, text=hscforce] {$mg$};
    \fill[hscink] (R) circle (1.8pt);
  \end{tikzpicture}
  \captionof{figure}{Forces on a rotor rider.}\label{ch2:fig:rotor}
\nextside
  \begin{itemize}
    \item \textbf{Rotor.} In a \term{rotor} ride a drum spins and the floor drops away. The
          riders stay on the wall (\cref{ch2:fig:rotor}).
    \item The wall's normal force is horizontal, towards the axis. It is the centripetal
          force: $F_\mathrm{N} = mv^2/r$.
    \item Static friction acts up the wall and balances the weight: $F_\mathrm{f} = mg$.
    \item Riders stay up only if $mg \le \mu F_\mathrm{N}$. This sets a slowest safe speed.
          Faster, the friction \emph{limit} grows, but the actual friction stays $mg$.
  \end{itemize}
\end{sidebyside}

\begin{formulacard}[notonsheet, unbreakable]{Friction limits: flat bend and rotor}
  \cardformula{v_\mathrm{max} = \sqrt{\mu g r} \qquad\qquad v_\mathrm{min} = \sqrt{\frac{gr}{\mu}}}
  \begin{symbols}
    \sym{$v_\mathrm{max}$}{fastest speed on a flat bend without skidding}{\si{m.s^{-1}}}
    \sym{$v_\mathrm{min}$}{slowest speed at which rotor riders do not slide down}{\si{m.s^{-1}}}
    \sym{$\mu$}{coefficient of friction (tyres on road; rider on wall)}{no unit}
    \sym{$g$}{gravitational acceleration, \qty{9.8}{m.s^{-2}}}{\si{m.s^{-2}}}
    \sym{$r$}{radius of the bend, or of the drum}{\si{m}}
  \end{symbols}
  \whentouse{a flat bend ($v_\mathrm{max}$); a rotor, floor taken away ($v_\mathrm{min}$).
    Mass makes no difference.}
  \derive{at the limiting speed, friction needed $=$ largest friction. Flat bend:
    $mv_\mathrm{max}^2/r = \mu mg$, so $v_\mathrm{max}^2 = \mu g r$. Rotor:
    $mg = \mu F_\mathrm{N} = \mu m v_\mathrm{min}^2/r$, so $v_\mathrm{min}^2 = gr/\mu$.}
  \trap{$\mu F_\mathrm{N}$ is the \emph{largest} friction, not the actual friction.}
\end{formulacard}

%% ==================================================================
\section{Key experiments}\label{ch2:sec:exp}
\covers{M5.IQ2.01}

\begin{experiment}{Investigating $F_\mathrm{c} = mv^2/r$ with a whirling stopper}
\aim{To investigate how $F_\mathrm{c}$ depends on speed (Part A), radius (Part B) and mass
  (Part C).}
\method{\begin{sidebyside}[0.46]
    \begin{itemize}
      \item Whirl the stopper in a horizontal circle (\cref{ch2:fig:apparatus}). The
            hanging mass stays still, so the tension provides $F_\mathrm{c} = Mg$.
      \item The paper clip keeps the radius $r = L$ fixed.
      \item Time 20 turns (less reaction-time error), three times, and average.
            $T = \text{time} \div 20$; $v = 2\pi L/T$.
    \end{itemize}
  \nextside
    \centering
    \begin{tikzpicture}[x=1cm, y=1cm]
    \draw[hscmuted, dashed, line width=0.7pt] (0,0.5) ellipse (3.0 and 0.32);
    \draw[object] (-0.12,-0.75) rectangle (0.12,0.66);
    \draw[hscink, line width=0.8pt] (0,0.66) -- (3.0,0.5);
    \draw[displacement, <->] (0.05,1.30) -- (3.0,1.14) node[pos=0.5, above, diagramlabel] {$L$};
    \fill[hscink] (3.0,0.5) ellipse (0.16 and 0.13);
    \node[diagramlabel, anchor=south, inner sep=1pt, fill=hscexp!4!white] at (3.0,0.68) {stopper, mass $m$};
    \draw[force] (2.85,0.51) -- (1.95,0.56) node[pos=0.5, below, inner sep=1pt, fill=hscexp!4!white, text=hscforce] {$F_\mathrm{T}$};
    \draw[hscink, line width=0.8pt] (0,-0.75) -- (0,-1.35);
    \draw[hscink, fill=hscmuted!40, line width=0.6pt] (-0.12,-0.98) rectangle (0.12,-0.88);
    \node[diagramlabel, anchor=west] at (0.25,-0.93) {paper clip (keeps $L$ fixed)};
    \node[diagramlabel, anchor=west] at (0.25,-0.42) {tube, held in the hand};
    \draw[object] (-0.33,-1.8) rectangle (0.33,-1.35);
    \node[diagramlabel] at (0,-1.575) {$M$};
    \draw[force] (0,-1.8) -- (0,-2.3) node[right, text=hscforce] {$Mg$};
    \node[diagramlabel, anchor=west] at (0.45,-1.575) {hanging mass};
    \end{tikzpicture}
    \captionof{figure}{Apparatus (side view).}\label{ch2:fig:apparatus}
  \end{sidebyside}
  \begin{itemize}
    \item \textbf{Part A:} change $M$ (so $F_\mathrm{c}$); keep $m$ and $L$ the same. Plot
          $F_\mathrm{c}$ against $v^2$. \textbf{Part B:} change $L$ (keep $M$, $m$); plot $v^2$
          against $r$. \textbf{Part C:} use 1, 2, 3, 4 identical stoppers (keep $M$, $L$); plot
          $v^2$ against $1/m$.
    \item Part A: \term{independent variable} (changed) $M$; \term{dependent variable}
          (measured) the time for 20 turns; \term{controlled variables} (kept the same) $m$,
          $L$, a horizontal circle.
  \end{itemize}}
\results{Part A: $F_\mathrm{c}$ against $v^2$ is a straight line through the \term{origin}
  (0, 0). The gradient of the \term{line of best fit} (the straight line closest to all the
  points) agrees with $m/r$ within 1\,\%. Parts B and C (expected from theory): $v^2$ against
  $r$ and $v^2$ against $1/m$ should also be straight lines through the origin.}
\conclusion{Part A: $F_\mathrm{c} \propto v^2$. This supports $F_\mathrm{c} = mv^2/r$. If the
  Part B and C graphs are straight, then at a fixed speed $F_\mathrm{c} \propto 1/r$ and
  $F_\mathrm{c} \propto m$. ($F_\mathrm{c}$ against $1/T^2$ is also straight; $F_\mathrm{c}$
  against $v$ is a curve.)}
\significance{It confirms $F_\mathrm{c} = mv^2/r$ with a measurable force (a hanging weight).
  Engineers use the same model for bends, banked tracks, rides and satellite orbits.}
\end{experiment}

%% ==================================================================
\begin{summarybox}[unbreakable]{Remember}
\begin{itemize}
  \item Constant speed in a circle is still acceleration, because the direction changes.
        $a_\mathrm{c} = v^2/r$ points to the centre. Also $v = 2\pi r/T$,
        $f = 1/T$ and $\omega = \Delta\theta/t$.
  \item $F_\mathrm{c} = mv^2/r$ is the net force towards the centre, given by a real
        force (tension, normal force, friction, gravity). There is no centrifugal force.
  \item Draw only real forces. Towards the centre they add up to $mv^2/r$; vertically
        (horizontal circle) they add up to zero.
  \item Conical pendulum or banked road (design speed): $\tan\theta = v^2/(rg)$.
        Flat bend: $v_\mathrm{max} = \sqrt{\mu g r}$. Rotor: $v_\mathrm{min} = \sqrt{gr/\mu}$.
  \item Whirling stopper: $F_\mathrm{c} = Mg$; $F_\mathrm{c}$ against $v^2$ is a straight line
        through the origin, with gradient $m/r$.
\end{itemize}
\end{summarybox}
}{}% Circular Motion: Forces and Real Situations
\IfFileExists{./ch3.tex}{\setcounter{chapter}{2}% ch3.tex -- Module 5 ESSENTIALS, Chapter 3: Circular Motion: Angles, Energy, Work and Torque
% Knowledge points only (no questions). Condensed from src/mod5/ch3.tex (same notation,
% same facts). Dot points: M5.IQ2.03 (omega = Delta theta / t), M5.IQ2.04, M5.IQ2.05.
% All diagrams are TikZ copied (some shrunk) from the full chapter: no figure files needed.
% Every number is taken from the full book and checked by
\chapter{Circular Motion: Angles, Energy, Work and Torque}\label{ch3:chap}
\covers{M5.IQ2.04, M5.IQ2.05}

\begin{summarybox}{In this chapter}
\begin{itemize}
  \item Angles in radians, the angular velocity $\omega$ and the speed $v = \omega r$.
  \item Why the centripetal force does no work; how work changes the total energy.
  \item Speeds and forces in a vertical circle.
  \item Torque, $\tau = r_\perp F = rF\sin\theta$, and how torque changes rotation.
\end{itemize}
\end{summarybox}

%% ==================================================================
\section{Angles and angular velocity}\label{ch3:sec:omega}
\covers{M5.IQ2.03}

\begin{bigidea}
Every point on a spinning object turns through the same angle in the same time, so it has
the same angular velocity $\omega$. Points further out move faster.
\end{bigidea}

\begin{keywords}
  \kw{Radian (rad)}{a unit for angles. One full turn is $2\pi$ rad, the same as \ang{360}.}
  \kw{Angular displacement ($\Delta\theta$)}{the angle turned through, measured at the centre.}
  \kw{Angular velocity ($\omega$)}{the angle turned each second. Unit: \si{rad.s^{-1}}.}
\end{keywords}

\begin{sidebyside}[0.33]
  \centering
  \begin{tikzpicture}[scale=0.95]
    \def\R{1.8}
    \fill[hscbrand!6] (0,0) circle (\R);
    \draw[hscmuted!70, line width=0.8pt] (0,0) circle (\R);
    \draw[guide] (0,0) -- (-30:\R);
    \draw[guide] (0,0) -- (40:\R);
    \draw[displacement] (-30:0.85) arc (-30:40:0.85);
    \draw[displacement] (-30:1.7) arc (-30:40:1.7);
    \draw[hscink, line width=0.6pt] (-30:0.35) arc (-30:40:0.35);
    \node[diagramlabel, font=\footnotesize] at (5:0.6) {$\Delta\theta$};
    \node[diagramlabel, text=hscmuted] at (5:1.25) {$s_\mathrm{A}$};
    \node[diagramlabel, text=hscmuted, anchor=west] at (5:1.8) {$s_\mathrm{B}$};
    \coordinate (A) at (40:0.85);
    \coordinate (B) at (40:1.7);
    \draw[velocity] (A) -- ++(130:0.55) node[left, text=hscvel, inner sep=1pt] {$v_\mathrm{A}$};
    \draw[velocity] (B) -- ++(130:1.1) node[left, text=hscvel, inner sep=1pt] {$v_\mathrm{B}$};
    \fill[object] (A) circle (3.5pt);
    \fill[object] (B) circle (3.5pt);
    \node[diagramlabel, anchor=south, inner sep=1pt] at ($(A)+(80:0.2)$) {A};
    \node[diagramlabel, anchor=south, inner sep=1pt] at ($(B)+(80:0.2)$) {B};
    \fill[hscink] (0,0) circle (1.5pt);
  \end{tikzpicture}
  \captionof{figure}{A and B turn through the same angle in the same time.}\label{ch3:fig:sameomega}
\nextside
  \begin{itemize}
    \item One turn $= 2\pi \text{ rad} = \ang{360}$, so $1 \text{ rad} \approx \ang{57.3}$.
    \item To radians: degrees $\times\dfrac{\pi}{180}$, turns $\times\, 2\pi$. \term{rpm} (turns per
          minute) $\times\dfrac{2\pi}{60}$ gives \si{rad.s^{-1}}.
    \item \term{Arc length} (distance along the edge): $s = r\theta$, with $\theta$ in rad (not on
          the sheet).
    \item In \cref{ch3:fig:sameomega}, A and B have the same $\omega$. B is twice as far out, so
          its arc is twice as long and it moves twice as fast.
    \item Every point has the same $\omega$, $T$ and $f$. The \term{tangential speed}
          $v = \omega r$ (speed along the circle) and $a_\mathrm{c} = \omega^2 r$ grow with $r$
          (same $\omega$).
  \end{itemize}
\end{sidebyside}

\begin{formulacard}{Angular velocity}
  \cardformula{\omega = \frac{\Delta\theta}{t}}
  \begin{symbols}
    \sym{$\omega$}{angular velocity: angle turned each second}{\si{rad.s^{-1}}}
    \sym{$\Delta\theta$}{angular displacement: angle turned}{\si{rad}}
    \sym{$t$}{time taken}{\si{s}}
  \end{symbols}
  \whentouse{anything that spins or moves round a circle: wheels, fans.}
  \trap{one turn is $2\pi$ rad, not 1 rad. Change turns or degrees to radians.}
\end{formulacard}

\textbf{From the period, and the speed of a point} \hscbadgenotonsheet
\[ \omega = \frac{2\pi}{T} = 2\pi f \qquad\qquad v = \omega r \qquad\qquad a_\mathrm{c} = \omega^2 r \]
\begin{itemize}
  \item $T$ = period (s), $f$ = frequency (Hz), $r$ = radius (m), not the diameter.
  \item Derive: one turn ($2\pi$ rad) takes one period, so $\omega = \dfrac{2\pi}{T}$. The
        sheet gives $v = \dfrac{2\pi r}{T}$, so ${v = \dfrac{2\pi}{T}\times r = \omega r}$. Put
        $v = \omega r$ into $a_\mathrm{c} = \dfrac{v^2}{r}$: $a_\mathrm{c} = \omega^2 r$.
\end{itemize}

%% ==================================================================
\section{Work and energy in circular motion}\label{ch3:sec:work}
\covers{M5.IQ2.04}

\begin{bigidea}
In uniform circular motion the net force is at \ang{90} to the motion. So it does no work,
and the speed stays constant.
\end{bigidea}

\begin{keywords}
  \kw{Work ($W$)}{energy transferred by a force when the object moves along the direction of
    the force (or of part of it). Unit: joule (J).}
  \kw{Total mechanical energy ($E$)}{kinetic energy plus potential energy: $E = K + U$.}
\end{keywords}

\begin{sidebyside}[0.27]
  \centering
  \begin{tikzpicture}[x=1cm, y=1cm]
    \def\R{1.1}
    \draw[hscmuted!60, line width=0.8pt] (0,0) circle (\R);
    \fill[hscink] (0,0) circle (1.5pt);
    \node[diagramlabel, left, inner sep=3pt, font=\footnotesize] at (0,0) {centre};
    \coordinate (P) at (\R,0);
    \draw[hscink, line width=0.5pt] ($(P)+(-0.22,0)$) -- ($(P)+(-0.22,0.22)$) -- ($(P)+(0,0.22)$);
    \draw[force] (P) -- ++(-0.75,0) node[pos=0.5, below, inner sep=3pt, text=hscforce] {$F_\mathrm{c}$};
    \draw[velocity] (P) -- ++(0,0.95) node[right, inner sep=2pt, text=hscvel] {$v$};
    \draw[displacement] ($(P)+(0.3,0.0)$) -- ++(0,0.45) node[right, inner sep=2pt, text=hscmuted, font=\footnotesize] {$\Delta s$};
    \fill[object] (P) circle (4.5pt);
  \end{tikzpicture}
  \captionof{figure}{The net force is at \ang{90} to the velocity: no work.}\label{ch3:fig:work}
\nextside
  \begin{itemize}
    \item $W = F_\parallel s = Fs\cos\theta$ \hscbadgeonsheet, with $\theta$ the angle between
          the force and the displacement $s$.
    \item The \term{centripetal force} (the net force towards the centre) is at \ang{90} to the
          velocity at every instant (\cref{ch3:fig:work}). So for each tiny step,
          $W = F_\mathrm{c}\,\Delta s\cos\ang{90} = 0$.
    \item \term{Work--energy theorem}: the work done by the net force equals the change in
          kinetic energy, $W_\mathrm{net} = \Delta K$ (not on the sheet). So
          $\Delta K = 0$: $K = \tfrac{1}{2}mv^2$ and the speed stay constant. Only the
          \emph{direction} of the velocity changes.
  \end{itemize}
\end{sidebyside}

\begin{itemize}
  \item If it speeds up, a \emph{tangential} part of the net force (along the path) does work.
  \item \textbf{Horizontal circle or circular orbit:} the height (or orbit radius) does not
        change, so $U$ is constant. $K$ is constant too, so $E = K + U$ is constant.
  \item \textbf{Ferris wheel:} steady speed, so $K$ stays the same, but $U$ grows as a rider
        rises. The seat does positive work on the rider. The energy comes from the motor.
\end{itemize}

\begin{formulacard}[notonsheet]{Work and total energy}
  \cardformula{W_\mathrm{other} = \Delta E = \Delta K + \Delta U}
  \begin{symbols}
    \sym{$W_\mathrm{other}$}{work done by all forces except gravity}{\si{J}}
    \sym{$\Delta E$}{change in total mechanical energy}{\si{J}}
    \sym{$\Delta K$, $\Delta U$}{change in kinetic and in potential energy}{\si{J}}
  \end{symbols}
  \whentouse{a push, a motor or friction does work on the object.}
  \derive{gravity's work is $-\Delta U$, so $W_\mathrm{net} = -\Delta U + W_\mathrm{other} = \Delta K$.
    Add $\Delta U$ to both sides.}
  \trap{do not count gravity twice. Its work is already in $\Delta U$.}
\end{formulacard}

\subsection{Vertical circles}

If only gravity does work (no motor, push or friction), the speed changes round a vertical
circle: \term{non-uniform circular motion} (moving in a circle with a changing speed).

\begin{sidebyside}[0.4]
  \centering
  \begin{tikzpicture}[scale=0.8]
    \def\R{1.7}
    \draw[hscmuted!60, line width=0.8pt] (0,0) circle (\R);
    \draw[hscink!60, line width=0.6pt] (0,0) -- (0,\R);
    \draw[hscink!60, line width=0.6pt] (0,0) -- (\R,0);
    \draw[hscink!60, line width=0.6pt] (0,0) -- (0,-\R);
    \fill[hscink] (0,0) circle (1.5pt);
    \node[diagramlabel, left, inner sep=3pt, font=\footnotesize] at (0,0) {centre};
    % top
    \draw[force] (-0.11,\R) -- (-0.11,1.1) node[left, inner sep=1pt, text=hscforce] {$F_\mathrm{T}$};
    \draw[force] (0.11,\R) -- (0.11,0.85) node[right, inner sep=1pt, text=hscforce] {$mg$};
    \draw[velocity] (0,\R) -- (-0.7,\R) node[left, inner sep=2pt, text=hscvel] {$v$};
    \fill[object] (0,\R) circle (4.5pt);
    \node[diagramlabel, above, inner sep=6pt] at (0,\R) {slowest};
    % side
    \draw[force] (\R,0) -- (0.9,0) node[above, inner sep=2pt, text=hscforce] {$F_\mathrm{T}$};
    \draw[force] (\R,0) -- (\R,-0.85) node[right, inner sep=2pt, text=hscforce] {$mg$};
    \draw[velocity] (\R,0) -- (\R,1.05) node[right, inner sep=2pt, text=hscvel] {$v$};
    \fill[object] (\R,0) circle (4.5pt);
    % bottom
    \draw[force] (0,-\R) -- (0,-0.4) node[right, pos=0.6, inner sep=2pt, text=hscforce] {$F_\mathrm{T}$};
    \draw[force] (0,-\R) -- (0,-2.55) node[right, inner sep=2pt, text=hscforce] {$mg$};
    \draw[velocity] (0,-\R) -- (1.5,-\R) node[right, inner sep=2pt, text=hscvel] {$v$};
    \fill[object] (0,-\R) circle (4.5pt);
    \node[diagramlabel, anchor=north east, inner sep=3pt] at (-0.15,-\R) {fastest};
  \end{tikzpicture}
  \captionof{figure}{A ball on a string in a vertical circle (arrows not to scale).}\label{ch3:fig:vertical}
\nextside
  \begin{itemize}
    \item The tension $F_\mathrm{T}$ points along the string, to the centre, at \ang{90} to the
          motion (\cref{ch3:fig:vertical}). It does no work.
    \item Gravity points down. Going up, part of it slows the ball. Going down, it speeds the
          ball up. So the ball is slowest at the top and fastest at the bottom.
    \item $K + U$ stays the same: the \term{conservation of mechanical energy}.
  \end{itemize}
\end{sidebyside}

\par\noindent\begin{minipage}{\linewidth}
\textbf{Speed: use energy} (no friction):
\begin{center}
  $\displaystyle \tfrac{1}{2}mv_1^2 + mgh_1 = \tfrac{1}{2}mv_2^2 + mgh_2$\\[4pt]
  \hscbadgecombine{$K = \frac{1}{2}mv^2$ and $\Delta U = mg\Delta h$}
\end{center}
Multiply by $\dfrac{2}{m}$: $v_2^2 = v_1^2 + 2g(h_1 - h_2)$. Top to bottom, the drop is the
diameter, $2r$ (not $r$), so $v_\mathrm{bottom}^2 = v_\mathrm{top}^2 + 4gr$.
\end{minipage}
\smallskip

\par\noindent\begin{minipage}{\linewidth}
\textbf{Forces: only at the very top or bottom.} Add the real forces, taking ``towards the
centre'' as positive. Their sum is $F_\mathrm{c} = \dfrac{mv^2}{r}$ (on the sheet), not an extra force.
\begin{center}
  $\displaystyle \text{top: } F_\mathrm{T} + mg = \frac{mv^2}{r} \qquad\qquad
    \text{bottom: } F_\mathrm{T} - mg = \frac{mv^2}{r}$\\[4pt]
  \hscbadgenotonsheet
\end{center}
\end{minipage}
\begin{itemize}
  \item At the top, $mg$ points \emph{to} the centre: $+mg$, not $-mg$.
  \item On a track, the \term{normal force} $F_\mathrm{N}$ (the push of the track or seat) replaces
        $F_\mathrm{T}$. In a \term{dip} (low point of a valley): $F_\mathrm{N} - mg = \dfrac{mv^2}{r}$.
        On a \term{crest} (top of a hill): $mg - F_\mathrm{N} = \dfrac{mv^2}{r}$.
  \item The string goes \emph{slack} (loose) at the top when $F_\mathrm{T} = 0$. Then
        $mg = \dfrac{mv^2}{r}$, so $v = \sqrt{gr}$: the smallest speed at the top. A cart leaves a
        crest if $v > \sqrt{gr}$.
  \item \term{Apparent weight} is the seat's push $F_\mathrm{N}$: above $mg$ in a dip, below $mg$
        on a crest.
\end{itemize}

%% ==================================================================
\section{Torque: what makes things turn}\label{ch3:sec:torque}
\covers{M5.IQ2.05}

\begin{bigidea}
A force can make an object turn about a pivot. Its turning effect is the torque,
$\tau = r_\perp F$. If the torques balance, the rotation does not change.
\end{bigidea}

\begin{keywords}
  \kw{Torque ($\tau$)}{the turning effect of a force. Unit: newton metre (\si{N.m}).}
  \kw{Pivot (axis)}{the point or line that an object turns about, such as a hinge.}
  \kw{Line of action}{the line along which a force acts, drawn longer in both directions.}
  \kw{Perpendicular distance ($r_\perp$)}{the shortest distance, pivot to line of action.}
\end{keywords}

\begin{itemize}
  \item Torque is bigger when the force is bigger, acts further from the pivot (a door handle
        is far from the hinges), and is closer to \ang{90} to the \term{lever} (the bar that turns).
  \item \Cref{ch3:fig:torque} shows two ways to find the torque (same result). $\theta$ is between
        $r$ and $F$. $\sin\ang{90} = 1$ gives the most torque. A force through the pivot gives none.
\end{itemize}

\begin{figure}
  \centering
  \begin{tikzpicture}[x=1cm, y=1cm]
    % ---- (a) perpendicular distance
    \begin{scope}
      \draw[object, rounded corners=1.5pt] (0.25,-0.13) rectangle (3.25,0.13);
      \draw[object] (0:0.34) -- (60:0.34) -- (120:0.34) -- (180:0.34) -- (240:0.34) -- (300:0.34) -- cycle;
      \fill[hscink] (0,0) circle (1.6pt);
      \node[diagramlabel, below, inner sep=5pt] at (0,-0.3) {pivot};
      \coordinate (P) at (3.0,0);
      \coordinate (Q) at (2.25,-1.299);
      \draw[guide] (P) -- (3.9,0);
      \draw[guide] ($(Q)+(240:0.35)$) -- (P);
      \draw[hscbrand, line width=1.1pt] (0,0) -- node[below left, diagramlabel, text=hscbrand, inner sep=1pt] {$r_\perp$} (Q);
      \draw[hscink, line width=0.5pt] ($(Q)+(150:0.2)$) -- ($(Q)+(150:0.2)+(60:0.2)$) -- ($(Q)+(60:0.2)$);
      \draw[hscink, <->, line width=0.5pt] (0,0.42) -- node[above, diagramlabel] {$r$} (3.0,0.42);
      \draw[hscink, line width=0.6pt] ($(P)+(0:0.5)$) arc (0:60:0.5);
      \node[diagramlabel] at ($(P)+(30:0.75)$) {$\theta$};
      \draw[force] (P) -- ++(60:1.3) node[above, inner sep=2pt, text=hscforce] {$F$};
      \fill[hscink] (P) circle (1.4pt);
      \node[diagramlabel] at (1.6,-1.95) {\textbf{(a)} $r_\perp = r\sin\theta$, \ $\tau = r_\perp F$};
    \end{scope}
    % ---- (b) perpendicular force
    \begin{scope}[xshift=6.9cm]
      \draw[object, rounded corners=1.5pt] (0.25,-0.13) rectangle (3.25,0.13);
      \draw[object] (0:0.34) -- (60:0.34) -- (120:0.34) -- (180:0.34) -- (240:0.34) -- (300:0.34) -- cycle;
      \fill[hscink] (0,0) circle (1.6pt);
      \node[diagramlabel, below, inner sep=5pt] at (0,-0.3) {pivot};
      \draw[hscink, <->, line width=0.5pt] (0,0.42) -- node[above, diagramlabel] {$r$} (2.92,0.42);
      \coordinate (P) at (3.0,0);
      \draw[guide] ($(P)+(0.65,0)$) -- ($(P)+(60:1.3)$) -- ($(P)+(0,1.126)$);
      \draw[force, hscforce!40, line width=1pt] (P) -- ++(60:1.3) node[right, inner sep=2pt, text=hscforce!70] {$F$};
      \draw[force] (P) -- ++(0,1.126) node[left, inner sep=5pt, text=hscforce, align=right] {$F\sin\theta$\\(turns)};
      \draw[force] (P) -- ++(0.65,0) node[right, inner sep=3pt, text=hscforce, align=left] {$F\cos\theta$\\(no turning)};
      \draw[hscink, line width=0.6pt] ($(P)+(0:0.35)$) arc (0:60:0.35);
      \node[diagramlabel] at ($(P)+(25:0.58)$) {$\theta$};
      \fill[hscink] (P) circle (1.4pt);
      \node[diagramlabel] at (1.6,-1.95) {\textbf{(b)} $\tau = r \times F\sin\theta$};
    \end{scope}
  \end{tikzpicture}
  \caption{A force $F$ on a spanner: (a) perpendicular distance; (b) perpendicular force.}\label{ch3:fig:torque}
\end{figure}

\begin{formulacard}{Torque}
  \cardformula{\tau = r_\perp F = rF\sin\theta}
  \begin{symbols}
    \sym{$\tau$}{torque: turning effect of the force}{\si{N.m}}
    \sym{$r_\perp$}{perpendicular distance, pivot to line of action}{\si{m}}
    \sym{$F$}{size of the force}{\si{N}}
    \sym{$r$}{distance from the pivot to where the force acts}{\si{m}}
    \sym{$\theta$}{angle between $r$ and $F$ (degree mode)}{\si{\degree}}
  \end{symbols}
  \whentouse{a force turns an object about a pivot: doors, spanners, motors.}
  \trap{measure $r$ from the pivot. If the angle is given from a line at \ang{90} to the lever,
    use $\theta = \ang{90}$ minus that angle. The unit is \si{N.m}, not J.}
\end{formulacard}

\subsection{How torque changes rotation}

A force changes straight-line motion. In the same way, a torque changes turning motion.
\begin{itemize}
  \item Net torque not zero: the object starts to turn, speeds up or slows down. A bigger net
        torque changes the rotation faster.
  \item \term{Balanced} torques (net torque zero): the rotation does not change. A still object
        stays still, and a turning object keeps turning at the same rate. About one pivot:
        $\sum\tau_\mathrm{clockwise} = \sum\tau_\mathrm{anticlockwise}$ \hscbadgenotonsheet
  \item Uses: door handles far from the hinges, long spanners. A bicycle \term{crank} (pedal arm)
        gives the most torque when it is horizontal and the foot pushes straight down.
\end{itemize}

%% ==================================================================
\section{Key experiments}\label{ch3:sec:experiments}
\covers{M5.IQ2.04, M5.IQ2.05}

\begin{experiment}{Investigating torque: $\tau = r_\perp F$}
\aim{To test $\tau = rF\sin\theta$, and to see how an unbalanced torque changes the rotation.}
\method{A metre rule is pivoted at its \term{centre of mass} (where its weight seems to act). A
  spring balance keeps it level against a fixed load, at different $r$ and angles. Then the
  torques are made unbalanced.}
\results{$Fr$ stayed at about \qty{0.59}{N.m}: $F$ against $1/r$ is a straight line through the
  origin. At an angle, $F\sin\theta$ stayed the same. A bigger net torque turned the rule faster.}
\conclusion{Only the part of $F$ at \ang{90} to the rule turns it: $\tau = rF\sin\theta$.}
\significance{It explains door handles, spanners and levers, and motor torque (Module 6).}
\end{experiment}

\begin{experiment}{Work and energy in circular motion}
\aim{To relate the work done on an object moving in a circle to its total energy.}
\method{\textbf{A:} a stopper on a string is whirled in a horizontal circle. A video gives
  $K = \tfrac{1}{2}mv^2$ at 8 points. \textbf{B:} a ball swings down from rest through a height
  $h$. A \term{light gate} (a beam that times how long the ball blocks it) gives its speed.}
\results{\textbf{A:} $K$ stayed at \qty{0.24}{J} all the way round, at the same height.
  \textbf{B:} $K$ at the bottom was 97\,\% of the $U$ lost ($mgh$), at every height.}
\conclusion{\textbf{A:} $E = K + U$ was constant: the tension is at \ang{90} to the motion and
  does no work. \textbf{B:} gravity does the work. The missing 3\,\% is a \term{systematic
  error} (the same in every reading): the light-gate beam is not a thin line, so every speed
  comes out slightly too low. Within this error, the $K$ gained equals the $U$ lost.}
\significance{Energy conservation explains pendulums, rollercoasters and satellites.}
\end{experiment}

%% ==================================================================
\begin{summarybox}{Remember}
\begin{itemize}
  \item One turn is $2\pi$ rad. $\omega = \Delta\theta/t = 2\pi/T$, $v = \omega r$:
        every point has the same $\omega$.
  \item Uniform circular motion: $F_\mathrm{c}$ is at \ang{90} to $v$, so $W = 0$ and $K$ is constant.
  \item Only forces other than gravity change $E$: $W_\mathrm{other} = \Delta E = \Delta K + \Delta U$.
  \item Vertical circle: energy for the speed (top to bottom: drop $2r$). Top:
        $F_\mathrm{T} + mg = mv^2/r$; bottom: $F_\mathrm{T} - mg = mv^2/r$.
  \item $\tau = r_\perp F = rF\sin\theta$, $r$ from the pivot. A net torque changes the rotation.
\end{itemize}
\end{summarybox}
}{}% Circular Motion: Angles, Energy, Work and Torque
\IfFileExists{./ch4.tex}{\setcounter{chapter}{3}% ch4.tex -- Module 5 ESSENTIALS, Chapter 4: Gravity, orbits and Kepler's laws
% Condensed from src/mod5/ch4.tex (knowledge points only, no questions).
% Dot points: M5.IQ3.01, M5.IQ3.02, M5.IQ3.03, M5.IQ3.04
% Numbers: the full book's verified values; re-checked by
% Figures: figures/ch4_cannonball.pdf, ch4_equalareas.pdf, ch4_ellipseforces.pdf
%          (copied from mod5/figures with their scripts)
\chapter{Gravity, Orbits and Kepler's Laws}\label{ch4:chap}
\covers{M5.IQ3.01, M5.IQ3.02, M5.IQ3.03, M5.IQ3.04}

\begin{summarybox}{In this chapter}
\begin{itemize}
  \item Newton's law of universal gravitation and the field strength $g$.
  \item Why a satellite stays in orbit, and its speed, period and radius.
  \item Kepler's three laws, explained with forces and energy.
  \item Low Earth orbits and geostationary orbits, and their uses.
\end{itemize}
\end{summarybox}

%% ==================================================================
\section{Newton's law of gravitation and field strength}\label{ch4:sec:newton}
\covers{M5.IQ3.01}

\begin{bigidea}
All masses attract each other. The pull weakens quickly with distance.
\end{bigidea}

\begin{keywords}
  \kw{Gravitational force}{the pull between any two masses. Unit: newton (N).}
  \kw{Inverse-square law}{double the distance, and the effect is four times smaller.}
  \kw{Field strength ($g$)}{the gravitational force on each kilogram at a point.}
\end{keywords}

\begin{formulacard}{Newton's law of universal gravitation}
  \cardformula{F = \frac{GMm}{r^2}}
  \begin{symbols}
    \sym{$F$}{force on each object (equal and opposite)}{\si{N}}
    \sym{$G$}{universal gravitational constant, \num{6.67e-11}}{\si{N.m^2.kg^{-2}}}
    \sym{$M$, $m$}{masses of the two objects}{\si{kg}}
    \sym{$r$}{distance between the two centres of mass}{\si{m}}
  \end{symbols}
  \whentouse{any two masses (the law is \term{universal}: it works everywhere), treated as
    points or spheres, such as planets and moons.}
  \trap{$r$ is centre to centre, not surface to surface. Square all of $r$.}
\end{formulacard}

\subsection{Gravitational field strength}
\begin{itemize}
  \item A \term{gravitational field} is a region where any mass feels a gravitational force.
        $g = F/m$, in \si{N.kg^{-1}} = \si{m.s^{-2}}. In \term{free fall} (gravity is the only
        force), $g$ is the acceleration.
  \item A planet's field is \term{radial}: field lines point to the centre and spread out, so the
        field weakens with distance. Near the surface it is nearly \term{uniform} (the same
        size and direction everywhere), $g = \qty{9.8}{m.s^{-2}}$.
  \item The \term{altitude} $h$ is the height above the surface. Gravity formulas need
        $r = r_\mathrm{E} + h$ from the centre (a definition, not on the sheet). Data sheet:
        $r_\mathrm{E} = \qty{6.371e6}{m}$, $M_\mathrm{E} = \qty{6.0e24}{kg}$.
\end{itemize}

\begin{formulacard}[notonsheet]{Gravitational field strength}
  \cardformula{g = \frac{GM}{r^2}}
  \begin{symbols}
    \sym{$g$}{field strength (free-fall acceleration)}{\si{N.kg^{-1}} = \si{m.s^{-2}}}
    \sym{$G$}{universal gravitational constant}{\si{N.m^2.kg^{-2}}}
    \sym{$M$}{mass of the planet (or moon, or star)}{\si{kg}}
    \sym{$r$}{distance from the planet's \emph{centre} to the point}{\si{m}}
  \end{symbols}
  \whentouse{at a planet's surface ($r$ = its radius) or above it ($r = r_\mathrm{E} + h$).}
  \derive{$mg = \dfrac{GMm}{r^2}$ (the same force, written two ways). Divide both sides by $m$.}
  \trap{never use the altitude $h$ alone as $r$. Change km to m.}
\end{formulacard}

\textbf{Factors that affect $g$.}
\begin{itemize}
  \item Only the planet's mass ($g \propto M$) and the distance from its centre ($g \propto 1/r^2$):
        $r \times 2 \Rightarrow g \div 4$. The falling object's mass cancels, so all objects fall
        with the same acceleration (if air resistance is negligible).
  \item Data-sheet values give $GM_\mathrm{E} = \qty{4.0e14}{m^3.s^{-2}}$ and a surface $g$ of
        \qty{9.86}{m.s^{-2}}. The sheet's 9.8 differs only because $M_\mathrm{E}$ is rounded.
  \item Real Earth: $g \approx \qty{9.78}{m.s^{-2}}$ at the equator and \qty{9.83}{m.s^{-2}} at
        the poles (the equator bulges, and the spin reduces the measured $g$ there). Mountains and dense rock
        change $g$ slightly.
  \item $g$ is a vector towards the mass. Add the fields of two bodies as vectors.
\end{itemize}

%% ==================================================================
\section{Satellites in circular orbits}\label{ch4:sec:orbits}
\covers{M5.IQ3.02, M5.IQ3.04}

\begin{bigidea}
A satellite is always falling, but it moves sideways so fast that the ground curves away
beneath it. Gravity provides the centripetal force.
\end{bigidea}

\begin{sidebyside}[0.46]
  \centering
  \includegraphics{ch4_cannonball.pdf}
  \captionof{figure}{Newton's cannonball (mountain far too tall). C is a circular orbit:
    gravity $F$ is at \ang{90} to $v$.}\label{ch4:fig:cannon}
\nextside
  Newton imagined a cannon on a very tall mountain, with no air (\cref{ch4:fig:cannon}).
  \begin{itemize}
    \item A slow ball falls to the ground (A). A faster one lands further away (B).
    \item At about \qty{7.9}{km.s^{-1}} near the surface, the ground curves away as fast as the
          ball falls. It never lands: it is in \term{orbit} (C).
    \item Faster still, the orbit is an oval (D). Above \qty{11.2}{km.s^{-1}} the ball escapes.
    \item An orbit is projectile motion that never ends: free fall. Gravity points to the
          centre, so \emph{gravity provides the centripetal force}: $F = F_\mathrm{c}$.
  \end{itemize}
\end{sidebyside}

\begin{formulacard}[notonsheet]{Orbital speed (circular orbit)}
  \cardformula{v = \sqrt{\frac{GM}{r}}}
  \begin{symbols}
    \sym{$v$}{orbital speed (the speed for a circular orbit)}{\si{m.s^{-1}}}
    \sym{$G$}{universal gravitational constant}{\si{N.m^2.kg^{-2}}}
    \sym{$M$}{mass of the central body (\emph{not} the satellite)}{\si{kg}}
    \sym{$r$}{orbital radius, from the centre of the central body}{\si{m}}
  \end{symbols}
  \whentouse{any circular orbit around a planet, moon or star.}
  \derive{$\dfrac{GMm}{r^2} = \dfrac{mv^2}{r}$ (both on the sheet). Divide by $m$, multiply by
    $r$: $v^2 = \dfrac{GM}{r}$. Take the square root.}
  \trap{$M$ is the planet's mass, not the satellite's. $r$ is not the altitude.}
\end{formulacard}

\begin{itemize}
  \item The satellite's mass cancels: all objects at one height orbit at the same speed, with the
        same period (the orbit simulation in \cref{ch4:sec:experiments} shows this).
  \item A higher orbit is slower: $v \propto 1/\sqrt r$. The acceleration
        $a_\mathrm{c} = v^2/r = GM/r^2$ is the local $g$.
  \item Orbital velocity is along the tangent to the orbit, at \ang{90} to the radius.
  \item \term{Apparent weightlessness} (feeling no weight in free fall): on the ISS (\qty{420}{km}
        up) $g = \qty{8.7}{m.s^{-2}}$, but astronauts and station fall together, so the floor does
        not push: the \term{normal force} is zero.
\end{itemize}

\subsection{The period and Kepler's third law}
\begin{formulacard}{Kepler's third law (Newton's form)}
  \cardformula{\frac{r^3}{T^2} = \frac{GM}{4\pi^2}}
  \begin{symbols}
    \sym{$r$}{orbital radius (for an ellipse: the semi-major axis)}{\si{m}}
    \sym{$T$}{period: the time for one orbit}{\si{s}}
    \sym{$G$}{universal gravitational constant}{\si{N.m^2.kg^{-2}}}
    \sym{$M$}{mass of the central body}{\si{kg}}
  \end{symbols}
  \whentouse{to find $T$, $r$ (then $h = r - r_\mathrm{E}$) or the central mass $M$.}
  \derive{put $v = \dfrac{2\pi r}{T}$ (on the sheet) into $v^2 = \dfrac{GM}{r}$:
    $\dfrac{4\pi^2 r^2}{T^2} = \dfrac{GM}{r}$. Multiply by $r$, divide by $4\pi^2$.}
  \trap{$T$ in seconds. Use a cube root (not a square root) to find $r$.}
\end{formulacard}

%% ==================================================================
\section{Kepler's laws of planetary motion}\label{ch4:sec:kepler}
\covers{M5.IQ3.04}

\begin{bigidea}
Kepler's laws came from data. Newton's law of gravitation explains all three.
\end{bigidea}

Kepler used Tycho Brahe's careful data. He published his first two laws in 1609 and the third
in 1619. They are \term{empirical laws}: they describe but do not explain. Newton (1687) showed that a force with
$F \propto 1/r^2$ explains all three.

\textbf{First law (ellipses).} \emph{Each planet moves in an ellipse, with the Sun at one focus.}
An \term{ellipse} is an oval with two \term{foci} (fixed points on its long axis). The other focus
is empty. Planet orbits are nearly circles (the Earth's \term{eccentricity}, its stretch,
is only 0.017), so they are treated as circles. Comets' orbits are very stretched. The closest and farthest points from
the Sun are \term{perihelion} and \term{aphelion} (around the Earth: perigee, apogee). For an
ellipse, $r$ in Kepler's third law is the \term{semi-major axis} (half the longest diameter),
$a = (r_\mathrm{p} + r_\mathrm{a})/2$.

\begin{sidebyside}[0.56]
  \centering
  \includegraphics{ch4_equalareas.pdf}
  \captionof{figure}{Equal times, equal areas: $\Delta t_1 = \Delta t_2$, so $A_1 = A_2$ (to
    scale).}\label{ch4:fig:areas}
\nextside
  \textbf{Second law (equal areas).} \emph{A line from the Sun to a planet sweeps out equal
  areas in equal times} (\cref{ch4:fig:areas}). Near the Sun the line is short, so the planet
  covers a longer arc in the same time: it is \textbf{fastest at perihelion},
  \textbf{slowest at aphelion}.
\end{sidebyside}

\begin{sidebyside}[0.52]
  \centering
  \includegraphics{ch4_ellipseforces.pdf}
  \captionof{figure}{Velocity (blue): length shows the speed. Gravity (red): direction only.
    Dashed red: the part of gravity along the motion.}\label{ch4:fig:forces}
\nextside
  \textbf{With forces} (\cref{ch4:fig:forces}). A force does \term{work} (gives energy) only if
  part of it points along the motion.
  \begin{itemize}
    \item \textbf{Circle:} gravity is at \ang{90} to $v$. No work, so constant speed.
    \item \textbf{Moving closer:} part of gravity points along the motion, so the planet
          speeds up.
    \item \textbf{Moving away:} part of gravity points against the motion, so it slows down.
  \end{itemize}
\end{sidebyside}

\textbf{With energy.} Only gravity acts, so the \term{total energy} (kinetic $K$ plus potential
$U$) is constant. $U = -\dfrac{GMm}{r}$ (on the sheet) is lowest (most negative) at perihelion, so
$K = \tfrac12 mv^2$ is largest there. In a circular orbit $U$, $K$ and the speed are constant.
At perihelion and aphelion only, $r_\mathrm{p} v_\mathrm{p} = r_\mathrm{a} v_\mathrm{a}$ (from
equal areas; not on the sheet; distances from the Sun).

\textbf{Third law (periods).} \emph{The square of the period is proportional to the cube of the
orbital radius (semi-major axis)}: $T^2 \propto r^3$. The constant, $GM/4\pi^2$, depends only on
the central mass: all the Sun's planets share one value, and all the Earth's satellites share
another. For two orbits of one body, $r_1^3/T_1^2 = r_2^3/T_2^2$ in any units, such as
years and \term{astronomical units} (1~AU = average Earth--Sun distance, \qty{1.50e11}{m}).

%% ==================================================================
\section{Low Earth orbits and geostationary orbits}\label{ch4:sec:leogeo}
\covers{M5.IQ3.03}

\begin{bigidea}
The height of an orbit fixes its speed and period. Each height suits different jobs
(\cref{ch4:tab:leogeo}).
\end{bigidea}

\begin{itemize}
  \item A \term{low Earth orbit} (LEO) is \qtyrange{200}{2000}{km} up. Thin air causes
        \term{orbital decay} (slow loss of height). A \emph{polar} LEO passes over both poles.
        The Earth turns beneath it, so it passes over every part of the Earth.
  \item A \term{geostationary} satellite stays above one point on the equator. It needs:
        (1) a period of one day (strictly 23\,h\,56\,min); (2) a circular orbit above the
        equator; (3) motion west to east, like the Earth's spin. Kepler's third law then allows
        only one radius, \qty{4.23e7}{m} (6.6 Earth radii).
  \item A \term{geosynchronous} orbit has a one-day period but may be tilted. Then the satellite
        drifts north and south in the sky, so it is not geostationary.
\end{itemize}

\begin{table}
  \caption{Low Earth orbit and geostationary orbit compared.}\label{ch4:tab:leogeo}
  \centering
  \small
  \begin{tabular}{@{}L{25mm}L{57mm}L{58mm}@{}}
    \toprule
    \colhead{Property} & \colhead{Low Earth orbit} & \colhead{Geostationary orbit}\\
    \midrule
    Altitude & \qtyrange{200}{2000}{km} (ISS: \qty{420}{km}) & about \qty{36000}{km}\\
    Speed, period & about \qty{7.7}{km.s^{-1}}, about \qty{90}{min} & \qty{3.1}{km.s^{-1}}, 1 day\\
    In the sky & crosses the sky in minutes & fixed, so ground dishes never move\\
    View & small area, fine detail & about 40\,\% of the Earth; poor near the poles\\
    Signal delay & a few milliseconds & about \qty{0.24}{s} (up and down)\\
    Problems & orbital decay; needs many satellites & weaker signals; costly launch\\
    Uses & ISS, imaging, mapping, polar weather satellites, internet constellations
      & TV, communications, weather images of a whole hemisphere\\
    \bottomrule
  \end{tabular}
\end{table}

Each use comes from a property: a low orbit gives fine detail \emph{because} it is close to the
ground. In both orbits gravity provides the centripetal force.

%% ==================================================================
% layout QA (after the whole-book build): without this break the first box started at the foot
% of the page with only its Aim and the "Method" heading. Section 4.5 (three boxes) plus the
% Remember box fill exactly two pages.
\clearpage
\section{Key experiments}\label{ch4:sec:experiments}
\covers{M5.IQ3.01, M5.IQ3.02, M5.IQ3.04}

\begin{experiment}{Factors that affect the field strength \texorpdfstring{$g$}{g}}
\aim{To show that surface $g$ is proportional to $M/r^2$ ($M$: mass, $r$: radius of the body).}
\method{\term{Secondary data} (measured by others) for the Moon, Mercury, Mars, Venus and the
  Earth; $g$ plotted against $M/r^2$.}
\results{A straight line through the origin; gradient \qty{6.67e-11}{N.m^2.kg^{-2}}.}
\conclusion{$g \propto M/r^2$, and the gradient is $G$, as $g = GM/r^2$ predicts.}
\significance{$M = gr^2/G$ gives a planet's mass. Cavendish found $G$ in 1798 (``weighing the Earth'').}
\end{experiment}

\begin{experiment}{Kepler's third law and the moons of Jupiter}
\aim{To show that all of Jupiter's large moons share one value of $r^3/T^2$, and to find Jupiter's mass.}
\method{$r$ and $T$ (in seconds) of Io, Europa, Ganymede and Callisto; $r^3$ plotted against $T^2$.}
\results{A straight line through the origin; gradient \qty{3.205e15}{m^3.s^{-2}}.}
\conclusion{$r^3 \propto T^2$. Gradient $= GM/4\pi^2$, so Jupiter's mass is \qty{1.90e27}{kg}.}
\significance{Galileo (1610): the first objects seen orbiting another planet. This supported the
  Sun-centred model.}
\end{experiment}

% Orbit simulation: the full book's numbers (mod5/ch4.tex, "Orbits in a simulation");
\begin{experiment}{Orbits in a simulation: mass, radius, speed and period}
\aim{To find how a satellite's period and speed depend on its own mass and on its orbital
  radius.}
\method{A computer orbit model (planet mass \qty{6.0e24}{kg}). The satellite is launched at
  \ang{90} to the radius, at the speed that gives a circular path; 5 orbits are timed. Only the
  satellite's mass is changed (\qty{100}{kg}, \qty{1000}{kg}, \qty{5000}{kg}), then only $r$
  (\num{7.0e6} to \qty{4.0e7}{m}).}
\results{The same period for every mass ($T = \qty{5817}{s}$ at $r = \qty{7.0e6}{m}$). A bigger
  $r$ gave a longer $T$ and a lower $v$, but $v^2 r = \qty{4.00e14}{m^3.s^{-2}}$ at every radius.}
\conclusion{The period does not depend on the satellite's mass. $v^2 r = GM$, so
  $v \propto 1/\sqrt{r}$, as $v = \sqrt{GM/r}$ predicts.}
\significance{All objects at one height move together, whatever their mass: an astronaut floats
  beside the space station.}
\end{experiment}

%% ==================================================================
\begin{summarybox}{Remember}
\begin{itemize}
  \item $F = GMm/r^2$ (on the sheet), $g = GM/r^2$ (not on it). $r = r_\mathrm{E} + h$, in metres.
  \item Gravity provides the centripetal force: $v = \sqrt{GM/r}$ (not on the sheet).
  \item A higher orbit is slower. The satellite's mass cancels.
  \item Weightless astronauts fall with their craft: the normal force is zero, not gravity.
  \item $r^3/T^2 = GM/(4\pi^2)$ for all orbits around one body, with $T$ in seconds.
  \item Kepler: ellipse with the Sun at a focus; equal areas, so fastest at perihelion; $T^2 \propto r^3$.
  \item Elliptical orbit: gravity does work, but $K + U$ stays constant ($K$ largest at perihelion).
  \item LEO: low, about \qty{90}{min}, fine detail. Geostationary: \qty{36000}{km}, one day, above
        the equator, west to east.
\end{itemize}
\end{summarybox}
}{}% Gravity, Orbits and Kepler's Laws
\IfFileExists{./ch5.tex}{\setcounter{chapter}{4}% ch5.tex -- Module 5 ESSENTIALS, Chapter 5: Energy in gravitational fields
% Knowledge points only (no questions). Condensed from src/mod5/ch5.tex:
% same notation, same facts and numbers (all numbers come from the full book's
\chapter[Energy in Gravitational Fields]{Energy in Gravitational Fields: Orbits, Escape and Changing Orbits}\label{ch5:chap}
\covers{M5.IQ3.05}

\begin{summarybox}{In this chapter}
\begin{itemize}
  \item Why gravitational potential energy $U$ is negative.
  \item The energy of a satellite in circular and elliptical orbits (Kepler's laws).
  \item Escape velocity, changing orbit and orbital decay.
\end{itemize}
\end{summarybox}

%% ==================================================================
\section{Gravitational potential energy in a radial field}\label{ch5:sec:gpe}
\covers{M5.IQ3.05}

\begin{bigidea}
Far from the ground, gravity gets weaker, so $mg\Delta h$ works only near the
surface. With $U = 0$ very far away, $U = -\dfrac{GMm}{r}$ is negative
everywhere else.
\end{bigidea}

\begin{keywords}
  \kw{Gravitational potential energy ($U$)}{energy stored because of where an
    object is in a gravitational field. Unit: J.}
  \kw{Radial field}{a field whose lines all point straight at the planet's centre.}
  \kw{Infinity}{so far away that the planet's gravity is too small to matter.}
\end{keywords}

\textbf{Why $U$ is negative} (three linked points):
\begin{enumerate}
  \item We choose $U = 0$ at infinity, where the pull of gravity is zero.
  \item Moving away needs work against gravity, so $U$ increases as $r$ increases.
  \item $U$ reaches 0 only at infinity. So $U$ is negative at every real distance.
\end{enumerate}
The minus sign comes from our choice of zero. Only \emph{changes} in $U$ can be
measured.

\begin{itemize}
  \item Work done by gravity: $W_\text{by gravity} = -\Delta U$. \term{Work
        against gravity} (by an engine, at a steady speed) is all stored:
        $W_\text{against} = +\Delta U$.
  \item \textbf{Where $U$ comes from:} the work against gravity from $r_1$ to
        $r_2$ is the area under the force--distance graph of
        $F = \dfrac{GMm}{r^2}$:
        $\Delta U = GMm\left(\dfrac{1}{r_1} - \dfrac{1}{r_2}\right)$. With
        $r_1 = r$, $r_2 = \infty$ and $U(\infty) = 0$, this gives
        $U = -\dfrac{GMm}{r}$.
\end{itemize}

\begin{formulacard}{Gravitational potential energy}
  \cardformula{U = -\frac{GMm}{r}}
  \begin{symbols}
    \sym{$U$}{gravitational potential energy (zero at infinity)}{\si{J}}
    \sym{$G$}{universal gravitational constant, \num{6.67e-11}}{\si{N.m^2.kg^{-2}}}
    \sym{$M$, $m$}{mass of the planet or star, mass of the object}{\si{kg}}
    \sym{$r$}{distance from the \emph{centre} of the planet}{\si{m}}
  \end{symbols}
  \whentouse{at any distance from a planet, moon or star. For a change, subtract:
    $\Delta U = U_2 - U_1 = GMm\left(\dfrac{1}{r_1} - \dfrac{1}{r_2}\right)$.}
  \trap{$r = r_\mathrm{E} + \text{altitude}$. Keep the minus sign.}
\end{formulacard}

%% ==================================================================
\section{Kinetic, potential and total energy in orbit}\label{ch5:sec:orbit}
\covers{M5.IQ3.05, M5.IQ3.04}

\begin{bigidea}
In a circular orbit the total energy is $U + K = -\dfrac{GMm}{2r}$. A lower
satellite moves faster, but it has \emph{less} total energy.
\end{bigidea}

The \term{total mechanical energy} $E$ is kinetic plus potential energy:
$E = U + K$ (unit: J).

\begin{formulacard}[notonsheet]{Kinetic and total energy in a circular orbit}
  \cardformula{K = \frac{GMm}{2r} \qquad\qquad E = U + K = -\frac{GMm}{2r}}
  \begin{symbols}
    \sym{$K$}{kinetic energy of the satellite}{\si{J}}
    \sym{$E$}{total mechanical energy, $U + K$}{\si{J}}
    \sym{$G$}{universal gravitational constant, \num{6.67e-11}}{\si{N.m^2.kg^{-2}}}
    \sym{$M$, $m$}{mass of the planet, mass of the satellite}{\si{kg}}
    \sym{$r$}{radius of the orbit, from the centre of the planet}{\si{m}}
  \end{symbols}
  \whentouse{\emph{circular} orbits only. Shortcuts: $E = -K$ and $U = 2E$.}
  \derive{gravity is the centripetal force.
    \begin{align*}
      \frac{mv^2}{r} &= \frac{GMm}{r^2}
        && \text{$F_\mathrm{c} = F$, so $v^2 = GM/r$ and $v = \sqrt{GM/r}$}\\
      K = \tfrac{1}{2}mv^2 &= \frac{GMm}{2r} && \text{put in } v^2 = GM/r\\
      E = U + K &= -\frac{2GMm}{2r} + \frac{GMm}{2r} = -\frac{GMm}{2r}
        && \text{add } U = -GMm/r = -2GMm/(2r)
    \end{align*}}
  \trap{for an ellipse use $K = \tfrac{1}{2}mv^2$. ``$E$ increases'' means it
    gets \emph{less} negative.}
\end{formulacard}


\Cref{ch5:fig:KUE} and \cref{ch5:tab:trends} compare circular orbits of
different radius.

\begin{sidebyside}[0.48]
  \centering
  \begin{tikzpicture}
    \begin{axis}[hsc, width=0.96\linewidth, height=0.62\linewidth,
        xmin=1, xmax=6, ymin=-66, ymax=36, xtick={1,2,...,6},
        ytick={-60,-40,-20,0,20},
        xlabel={$r$ (in units of $r_\mathrm{E}$)},
        ylabel={energy of \qty{1}{kg} (\si{MJ})}]
      \draw[hscink!70, line width=0.8pt] (axis cs:1,0) -- (axis cs:6,0);
      \addplot[hscseriesA, domain=1:6, samples=80, line width=1.2pt, solid] {31.41/x};
      \addplot[hscseriesB, domain=1:6, samples=80, line width=1.2pt, dashed] {-62.82/x};
      \addplot[hscseriesC, domain=1:6, samples=80, line width=1.2pt, densely dotted] {-31.41/x};
      \node[font=\small, text=hscseriesA, anchor=south] at (axis cs:4.2,7.9) {$K$};
      % E label above its own curve (E = -19.6 MJ at x = 1.6), clear of the U curve below
      \node[font=\small, text=hscseriesC, anchor=south] at (axis cs:1.6,-17.2) {$E$};
      \node[font=\small, text=hscseriesB, anchor=north] at (axis cs:2.5,-26.6) {$U$};
      \draw[guide] (axis cs:2,-31.41) -- (axis cs:2,15.70);
      \fill[hscseriesA] (axis cs:2,15.70) circle (1.6pt);
      \fill[hscseriesC] (axis cs:2,-15.70) circle (1.6pt);
      \fill[hscseriesB] (axis cs:2,-31.41) circle (1.6pt);
    \end{axis}
  \end{tikzpicture}
  \captionof{figure}{$K$, $U$ and $E$ for \qty{1}{kg} in a circular orbit around
    Earth. $U$ rises towards 0; $E = -K$ and $U = 2E$.}\label{ch5:fig:KUE}
\nextside
  \captionof{table}{Moving to a bigger circular orbit.}\label{ch5:tab:trends}
  \centering\small
  \begin{tabular}{@{}L{31mm}L{38mm}@{}}
    \toprule
    \colhead{Quantity} & \colhead{When $r$ increases}\\
    \midrule
    speed $v$ & decreases\\
    force of gravity $F$ & decreases ($F \propto 1/r^2$)\\
    kinetic energy $K$ & decreases\\
    potential energy $U$ & increases (less negative)\\
    total energy $E$ & increases (less negative)\\
    period $T$ & increases\\
    \bottomrule
  \end{tabular}

  \raggedright\vspace{2mm}
  If $E < 0$, the object is \term{bound} (trapped by gravity): it orbits or
  falls back. If $E \ge 0$, it escapes.
\end{sidebyside}

\subsection{Energy in elliptical orbits (Kepler's laws)}
\covers{M5.IQ3.04, M5.IQ2.04}

\textbf{Kepler's laws:} (1) each planet moves in an ellipse, with the Sun at one
focus; (2) a line from the Sun to the planet sweeps out equal areas in equal
times; (3) $\dfrac{r^3}{T^2} = \dfrac{GM}{4\pi^2}$ (on the sheet).
\Cref{ch5:fig:ellipse} links them to energy.

\begin{sidebyside}[0.45]
  \centering
  \begin{tikzpicture}[x=0.68cm, y=0.68cm]
    \draw[hscmuted!70, line width=0.8pt] (0,0) ellipse (2.4 and 2.078);
    \draw[->, hscmuted, line width=0.6pt] (60:2.75 and 2.40) arc (60:110:2.75 and 2.40);
    \fill[hscacc] (1.2,0) circle (5pt);
    \node[diagramlabel, font=\footnotesize, below=5pt] at (1.2,0) {Sun};
    \draw[velocity] (2.4,0) -- (2.4,1.3) node[right, text=hscvel, font=\footnotesize] {$v$};
    \fill[object] (2.4,0) circle (3.5pt);
    \draw[velocity] (-2.4,0) -- (-2.4,-0.433) node[left, text=hscvel, font=\footnotesize] {$v$};
    \fill[object] (-2.4,0) circle (3.5pt);
    \coordinate (P) at (-1.2,-1.80);
    % speed at P (r = 3.0) from energy conservation: 0.447 x perihelion speed, so length 0.58
    \draw[velocity] (P) -- ++(0.520,-0.260) node[below, text=hscvel, font=\footnotesize] {$v$};
    \draw[force] (P) -- ++(0.80,0.60) node[above left, text=hscforce, inner sep=1pt, font=\footnotesize] {$F$};
    \fill[object] (P) circle (3.5pt);
    \node[diagramlabel, font=\footnotesize, anchor=west] at (2.45,-0.35) {perihelion};
    \node[diagramlabel, font=\footnotesize, anchor=east] at (-2.45,0.35) {aphelion};
  \end{tikzpicture}
  \captionof{figure}{A planet in an elliptical orbit (anticlockwise). On the way
    in, part of $F$ points along $v$, so gravity speeds the planet
    up.}\label{ch5:fig:ellipse}
\nextside
  \begin{itemize}
    \item \textbf{Circle:} gravity is at \ang{90} to $v$, so
          $W = Fs\cos\ang{90} = 0$. The speed stays the same.
    \item \textbf{Ellipse:} gravity is the only force, so $E = K + U$ stays the
          same. When $U$ goes down, $K$ goes up.
    \item \textbf{Perihelion} (closest point to the Sun): $U$ lowest, $K$
          largest, fastest. \textbf{Aphelion} (farthest): slowest. This agrees
          with Kepler's second law.
  \end{itemize}
\end{sidebyside}

At any point of any path where gravity is the only force,
$E = K + U = \tfrac{1}{2}mv^2 - \dfrac{GMm}{r}$ (combine $K$ and $U$ from the
sheet). Use $v$ and $r$ at the \emph{same} point.

%% ==================================================================
\section{Escape velocity}\label{ch5:sec:escape}
\covers{M5.IQ3.05}

\begin{bigidea}
The \term{escape velocity} $v_\mathrm{esc}$ is the smallest launch speed that
lets a projectile get away from a planet for ever. At exactly this speed, its
total energy is zero.
\end{bigidea}

A \term{projectile} gets all its speed at the start, then moves with only
gravity acting (no engine).

The total energy decides the path. \Cref{ch5:fig:cannon} shows Newton's cannon
(on a very tall mountain, above the air) firing sideways at five speeds.

\begin{sidebyside}[0.50]
  \centering
  \begin{tikzpicture}[scale=0.88]
    \fill[hscseriesA!12] (0,0) circle (1);
    \draw[hscseriesA!60, line width=0.8pt] (0,0) circle (1);
    \fill[hscmuted!45] (-0.20,0.98) -- (0,1.3) -- (0.20,0.98) -- cycle;
    \draw[hscmuted, line width=1.1pt, dashed, domain=0:44.62, samples=50, variable=\t]
      plot ({0.637/(1-0.51*cos(\t))*sin(\t)}, {0.637/(1-0.51*cos(\t))*cos(\t)});
    \draw[hscseriesA, line width=1.1pt] (0,0) circle (1.3);
    \draw[hscseriesC, line width=1.1pt, domain=0:360, samples=150, variable=\t]
      plot ({1.573/(1+0.21*cos(\t))*sin(\t)}, {1.573/(1+0.21*cos(\t))*cos(\t)});
    \draw[hscseriesB, line width=1.1pt, domain=0:108, samples=90, variable=\t]
      plot ({2.6/(1+cos(\t))*sin(\t)}, {2.6/(1+cos(\t))*cos(\t)});
    \draw[hscforce, line width=1.1pt, densely dashed, domain=0:96, samples=90, variable=\t]
      plot ({3.328/(1+1.56*cos(\t))*sin(\t)}, {3.328/(1+1.56*cos(\t))*cos(\t)});
    \fill[object] (0,1.3) circle (2.8pt);
    \node[diagramlabel, above, font=\footnotesize] at (0,1.36) {launch};
    \node[diagramlabel, font=\footnotesize] at (0,-0.25) {Earth};
    \draw[hscmuted, line width=0.4pt] (0.72,0.70) -- (0.55,0.30);
    \node[diagramlabel, font=\footnotesize, text=hscmuted] at (0.50,0.15) {1};
    \node[diagramlabel, font=\footnotesize, text=hscseriesA] at (-1.44,0) {2};
    \node[diagramlabel, font=\footnotesize, text=hscseriesC, anchor=north] at (0,-2.0) {3};
    \node[diagramlabel, font=\footnotesize, text=hscseriesB, anchor=north west] at (3.42,-1.05) {4};
    \node[diagramlabel, font=\footnotesize, text=hscforce, anchor=west] at (3.98,-0.42) {5};
  \end{tikzpicture}
  \captionof{figure}{Newton's cannon at five speeds (not to scale: the mountain
    is drawn far too tall).}\label{ch5:fig:cannon}
\nextside
  \begin{enumerate}
    \item slower than orbital speed: falls back
    \item $v = \sqrt{GM/r}$: circle
    \item between $\sqrt{GM/r}$ and $v_\mathrm{esc}$: ellipse
    \item $v = v_\mathrm{esc}$: just escapes
    \item $v > v_\mathrm{esc}$: escapes and is still moving when very far away
  \end{enumerate}
  Paths 1--3 have $E < 0$ (bound). Path 4 has $E = 0$: it reaches infinity
  with $K = 0$ and $U = 0$. Path 5 has $E > 0$.
\end{sidebyside}

\begin{formulacard}[notonsheet]{Escape velocity}
  \cardformula{v_\mathrm{esc} = \sqrt{\frac{2GM}{r}}}
  \begin{symbols}
    \sym{$v_\mathrm{esc}$}{escape velocity (really a speed: the direction does
      not matter)}{\si{m.s^{-1}}}
    \sym{$G$}{universal gravitational constant, \num{6.67e-11}}{\si{N.m^2.kg^{-2}}}
    \sym{$M$}{mass of the planet, moon or star}{\si{kg}}
    \sym{$r$}{distance from the centre where the object starts}{\si{m}}
  \end{symbols}
  \whentouse{a projectile, from the surface or from an orbit.}
  \derive{energy is conserved. In the smallest case the object reaches infinity
    with $K = 0$ and $U = 0$, so $K + U = 0$ at launch too.
    \begin{align*}
      \tfrac{1}{2}mv^2 - \frac{GMm}{r} = 0, \text{ so } \tfrac{1}{2}v^2 &= \frac{GM}{r}
        && \text{add $GMm/r$, divide by $m$ ($m$ cancels)}\\
      v_\mathrm{esc} &= \sqrt{\frac{2GM}{r}} && \text{multiply by 2, take the square root}
    \end{align*}}
  \trap{do not lose the 2. Without it you get the orbital speed $\sqrt{GM/r}$.}
\end{formulacard}



\begin{itemize}
  \item $m$ cancels, so a pebble and a spacecraft need the same escape speed (the
        heavier one needs more kinetic energy). From Earth's surface,
        $v_\mathrm{esc} = \qty{11.2}{km.s^{-1}}$.
  \item At the same $r$, $v_\mathrm{esc} = \sqrt{2}\times v_\text{orbital}$.
  \item A rocket need not reach $v_\mathrm{esc}$: its engines keep adding energy,
        and $v_\mathrm{esc}$ is smaller higher up. It escapes once $K + U$
        reaches zero.
\end{itemize}

%% ==================================================================
\section{Moving between orbits}\label{ch5:sec:transfer}
\covers{M5.IQ3.05, M5.IQ3.03}

\begin{bigidea}
To reach a higher orbit, the engines must add energy $\Delta E$. The satellite
ends up \emph{slower}: $K$ falls by $\Delta E$ and $U$ rises by $2\Delta E$. Air
drag does the opposite.
\end{bigidea}

A bigger orbit has a less negative $E$, so the engines supply
$\Delta E = E_2 - E_1$. Because $K = -E$ and $U = 2E$ in every circular orbit,
$\Delta K = -\Delta E$ and $\Delta U = 2\Delta E$.

\begin{formulacard}[notonsheet]{Energy to change circular orbit}
  \cardformula{\Delta E = \frac{GMm}{2}\left(\frac{1}{r_1} - \frac{1}{r_2}\right)}
  \begin{symbols}
    \sym{$\Delta E$}{change in total energy (smallest engine work)}{\si{J}}
    \sym{$G$}{universal gravitational constant, \num{6.67e-11}}{\si{N.m^2.kg^{-2}}}
    \sym{$M$, $m$}{mass of the planet, mass of the satellite}{\si{kg}}
    \sym{$r_1$, $r_2$}{radius of the starting orbit, radius of the new orbit}{\si{m}}
  \end{symbols}
  \whentouse{from one circular orbit to another (not a launch from the ground).}
  \derive{$\Delta E = E_2 - E_1 = -\dfrac{GMm}{2r_2} + \dfrac{GMm}{2r_1}$; take
    out $\dfrac{GMm}{2}$.}
  \trap{$\Delta E$ is not $\Delta U$. The engines supply only half of the rise in
    $U$; the other half comes from the fall in $K$.}
\end{formulacard}



% layout QA (whole-book build): slash fractions in this paragraph keep its lines evenly spaced, so
% Figure 5.4 fits on the same page as its text and the Key experiment box stays in one piece.
\textbf{From the ground.} On the pad $K \approx 0$, so
$E = U = -GMm/r_\mathrm{E}$. The smallest energy to reach a circular
orbit is $\Delta E = -GMm/(2r) + GMm/r_\mathrm{E}$. For
\qty{1}{kg}: \qty{32.6}{MJ} to a \qty{250}{km} orbit, \qty{58.1}{MJ} to
geostationary orbit (GEO: above the equator, one turn per day, used by TV and
weather satellites) and \qty{62.8}{MJ} to escape. So low Earth orbit (LEO: space
station, Starlink) is much cheaper.

\textbf{During a launch}, the fuel's chemical energy becomes extra $K$ and $U$
(and heat and sound). After the engines stop, the rocket coasts up: $K$ turns
into $U$, and $K + U$ stays constant.


\subsection{The Hohmann transfer}

A \term{Hohmann transfer} is a fuel-saving way to change orbit, using two short
engine firings (\term{burns}). \term{Perigee} and \term{apogee} are the closest
and farthest points of an orbit around Earth. \Cref{ch5:fig:hohmann} shows a
move from a \qty{250}{km} orbit to GEO.

\begin{sidebyside}[0.42]
  \centering
  \begin{tikzpicture}[x=0.95cm, y=0.95cm]
    \fill[hscseriesA!15] (0,0) circle (0.55);
    \draw[hscseriesA!60, line width=0.8pt] (0,0) circle (0.55);
    \node[diagramlabel, font=\footnotesize] at (0,0) {Earth};
    \draw[hscmuted, line width=0.9pt] (0,0) circle (1.0);
    \draw[hscmuted, line width=0.9pt] (0,0) circle (2.3);
    \draw[hscseriesC, line width=1.2pt, domain=0:180, samples=80, variable=\t]
      plot ({-0.65+1.65*cos(\t)}, {1.5166*sin(\t)});
    \draw[hscseriesC!60, line width=0.7pt, dotted, domain=180:360, samples=60, variable=\t]
      plot ({-0.65+1.65*cos(\t)}, {1.5166*sin(\t)});
    \draw[hscseriesC, line width=1.2pt, ->, domain=40:72, samples=20, variable=\t]
      plot ({-0.65+1.65*cos(\t)}, {1.5166*sin(\t)});
    \draw[->, hscmuted, line width=0.7pt] (100:1.2) arc (100:140:1.2);
    % radius guides start at the Earth's edge, so they do not cross the word "Earth"
    \draw[guide] (-60:0.55) -- (-60:1.0) node[pos=0.5, left, font=\footnotesize, text=hscmuted] {$r_1$};
    \draw[guide] (-125:0.55) -- (-125:2.3) node[pos=0.42, left, font=\footnotesize, text=hscmuted] {$r_2$};
    \draw[accel] (1.0,0) -- (1.0,0.85) node[right, text=hscacc, font=\footnotesize, align=left, fill=white, inner sep=1pt] {burn 1};
    \draw[accel] (-2.3,0) -- (-2.3,-0.85) node[left, text=hscacc, font=\footnotesize, align=right] {burn 2};
    \fill[object] (1.0,0) circle (3pt);
    \fill[object] (-2.3,0) circle (3pt);
    \node[diagramlabel, font=\footnotesize, anchor=north west, fill=white, inner sep=1pt] at (1.0,-0.05) {perigee};
    \node[diagramlabel, font=\footnotesize, anchor=south east] at (-2.32,0.05) {apogee};
    \node[diagramlabel, font=\footnotesize, text=hscseriesC, align=center] at (-0.4,1.85) {transfer orbit};
  \end{tikzpicture}
  \captionof{figure}{Hohmann transfer, low orbit ($r_1$) to GEO ($r_2$), not to
    scale.}\label{ch5:fig:hohmann}
\nextside
  \begin{enumerate}
    \item \textbf{Burn 1} at perigee: push \emph{forwards} (along the motion)
          into the elliptical \term{transfer orbit}.
    \item \textbf{Coast:} gravity slows it; $K$ turns into $U$, $E$ is constant.
    \item \textbf{Burn 2} at apogee: push forwards into the big circle.
  \end{enumerate}
  Both burns speed it up, but it ends \emph{slower} (7.77 to
  \qty{3.08}{km.s^{-1}}). To move \emph{down}, push backwards.
\end{sidebyside}

\subsection{Orbital decay}

\term{Orbital decay} is a slow fall to lower orbits. A LEO satellite moves
through very thin air. \term{Atmospheric drag} (the push of this air against its
motion) turns some mechanical energy into heat (\cref{ch5:fig:decay}).

\begin{sidebyside}[0.32]
  \centering
  \begin{tikzpicture}[x=0.8cm, y=0.8cm]
    \fill[hscseriesA!15] (0,0) circle (0.9);
    \draw[hscseriesA!60, line width=0.8pt] (0,0) circle (0.9);
    \node[diagramlabel, font=\footnotesize] at (0,-0.3) {Earth};
    \draw[hscmuted, line width=0.9pt, domain=0:1080, samples=240, variable=\t, ->]
      plot ({(2.0-0.6*\t/1080)*cos(\t)}, {(2.0-0.6*\t/1080)*sin(\t)});
    \coordinate (S) at ({1.975*cos(45)}, {1.975*sin(45)});
    \draw[velocity] (S) -- ++(135:0.95) node[above, text=hscvel, font=\footnotesize] {$v$};
    \draw[force] (S) -- ++(-45:0.55) node[right, text=hscforce, font=\footnotesize] {drag};
    % the "gravity" label sits inside the Earth disc, just below the arrow tip, so it covers
    % neither the Earth's outline nor the spiral
    \draw[force] (S) -- ($(S)!0.75cm!(0,0)$);
    \node[anchor=north east, text=hscforce, font=\footnotesize, inner sep=1pt] at (45:0.775) {gravity};
    \fill[object] (S) circle (3.5pt);
  \end{tikzpicture}
  \captionof{figure}{Drag makes the satellite spiral slowly inwards (spacing
    exaggerated).}\label{ch5:fig:decay}
\nextside
  \begin{enumerate}
    \item Drag removes energy, so $E$ falls (more negative).
    \item $E = -\dfrac{GMm}{2r}$, so $r$ falls. The orbit gets lower.
    \item $E = -K$, so $K$ \emph{rises}: the satellite speeds up.
    \item $U = 2E$, so $U$ falls by twice the energy lost: half becomes heat,
          half becomes extra $K$.
  \end{enumerate}
  Lower orbits meet denser air, so decay speeds up until the satellite burns
  up. Decay is caused by drag, not by ``stronger gravity''.
\end{sidebyside}

%% ==================================================================
\section{Key experiments}\label{ch5:sec:experiment}
\covers{M5.IQ3.05, M5.IQ3.04}

\begin{experiment}{Launch speed, total energy and escape (simulation)}
\aim{To find how the launch speed decides whether a probe falls back, orbits or
  escapes, and to test $v_\mathrm{esc} = \sqrt{2GM/r}$ using total energy.}
\method{A computer orbit model (Earth's mass) launches a probe sideways from
  $r_0 = \qty{1.00e7}{m}$ at 5 to \qty{9}{km.s^{-1}}. Each run gives the path
  and $E/m = \tfrac{1}{2}v_0^2 - GM/r_0$.}
\results{\begin{center}\small
  \begin{tabular}{@{}lccccc@{}}
    \toprule
    $v_0$ (\si{km.s^{-1}}) & 5.00 & 6.33 & 7.00 & 8.94 & 8.95\\
    $E/m$ (\si{MJ.kg^{-1}}) & $-27.5$ & $-20.0$ & $-15.5$ & $-0.06$ & $+0.03$\\
    path & falls back & circle & ellipse & very long ellipse & escapes\\
    \bottomrule
  \end{tabular}
  \end{center}}
\conclusion{The probe escapes only when $E \ge 0$, between 8.94 and
  \qty{8.95}{km.s^{-1}}; theory gives $\sqrt{2GM/r_0} = \qty{8.95}{km.s^{-1}}$.
  Around an ellipse, $E$ stays the same while $K$ and $U$ swap.}
\significance{Total energy, not speed alone, decides whether an object is bound.}
\end{experiment}

%% ==================================================================
\begin{summarybox}{Remember}
\begin{itemize}
  \item $U = -\dfrac{GMm}{r}$: zero at infinity, negative everywhere else, with
        $r = r_\mathrm{E} + \text{altitude}$.
  \item Circular orbit: $K = \dfrac{GMm}{2r}$ and
        $E = U + K = -\dfrac{GMm}{2r}$, so $E = -K$ and $U = 2E$.
  \item Gravity only: $E = K + U$ stays constant. An elliptical orbit is fastest
        at perihelion. $E < 0$: bound; $E \ge 0$: escapes.
  \item $K + U = 0$ gives $v_\mathrm{esc} = \sqrt{\dfrac{2GM}{r}}$, whatever the
        mass (\qty{11.2}{km.s^{-1}} from Earth's surface).
  \item A higher orbit needs
        $\Delta E = \dfrac{GMm}{2}\left(\dfrac{1}{r_1} - \dfrac{1}{r_2}\right)$:
        $E$ and $U$ rise, but $K$ falls.
  \item Orbital decay: drag lowers the orbit and speeds the satellite up.
\end{itemize}
\end{summarybox}
}{}% Energy in Gravitational Fields

%% ---------------------------------------------------------------- back
\IfFileExists{./formulas.tex}{% formulas.tex -- Module 5 ESSENTIALS: All formulas at a glance.
% One row per formula of the module, in chapter order. Formulas, notation and sheet status
% are copied from the full book's src/mod5/formulasheet.tex (which follows
% (F_f <= mu F_N, W_net = Delta K, v = sqrt(gr)) are results stated in the chapters
% (full book ch2/ch3 and these Essentials), all "not on the sheet".
% COMBINE), a_c = omega^2 r (full book ch3 card "Speed of a point") and the dip/crest normal force
% (full book and Essentials ch3) were added.
\unnumberedchapter[\faIcon{square-root-alt}]{All formulas at a glance}
\phantomsection\label{fg:chap}

% ---- helpers for this part only (prefix fg)
\newcommand\fgYes{\textcolor{hscformula}{\footnotesize\bfseries\faCheck~Yes}}
\newcommand\fgNo{\textcolor{hscmistake}{\footnotesize\bfseries\faTimes~Learn~it}}
\newcommand\fgComb{\textcolor{hsctiptext}{\footnotesize\bfseries\faLink~Combine}}
% a formula with 2.4pt of clear space above and below it, so fractions never touch the rules
\newcommand\fgf[1]{\raisebox{0pt}[\dimexpr\height+2.4pt\relax][\dimexpr\depth+2.4pt\relax]{#1}}
% \fgrow{mark}{formula}{what it is for}{chapter}
\newcommand\fgrow[4]{\fgf{#2} & #3 & #1 & #4\\\hline}
% \fgrowii{mark}{formula, line 1}{formula, line 2}{what it is for}{chapter}
\newcommand\fgrowii[5]{\fgf{#2}\newline\fgf{#3} & #4 & #1 & #5\\\hline}
% a darker rule between two chapters
\newcommand\fgnext{\noalign{\vskip-0.4pt{\color{hscbrand!45}\hrule height 0.9pt}}}
% chapter number, linked to the chapter
\newcommand\fgch[1]{\hyperref[ch#1:chap]{\textbf{#1}}}
\newcommand\fghead{\colhead{Formula} & \colhead{What it is for} & \colhead{Sheet} &
  \colhead{Ch.}}

Every formula in Module~\modulenumber, in chapter order, written as on the official formula
sheet. \fgYes: printed on the sheet, so you are given it in the exam. \fgNo: not on the
sheet, so learn it or learn how to derive it. \fgComb: not printed in this form; join
formulas from the sheet in one or two lines. ``Ch.'' is the chapter.

{\small
\setlength\LTpre{4pt}\setlength\LTpost{6pt}
\renewcommand\arraystretch{1.1}\arrayrulecolor{hscrule}
\begin{longtable}{@{}C{52mm}@{\hspace{3mm}}L{\dimexpr\linewidth-52mm-17mm-7mm-9mm\relax}%
    @{\hspace{3mm}}C{17mm}@{\hspace{3mm}}C{7mm}@{}}
  \fghead\\\hline\endfirsthead
  \fghead\\\hline\endhead
  % ---- Chapter 1: Projectile Motion
  \fgrow{\fgNo}{$u_x = u\cos\theta$ \quad $u_y = u\sin\theta$}
    {Parts of the launch velocity.}{\fgch{1}}
  \fgrowii{\fgYes}{$v = u + at$ \quad $v^2 = u^2 + 2as$}{$s = ut + \dfrac{1}{2}at^2$}
    {Constant acceleration, once in each direction: $a_x = 0$,
     $a_y = \qty{-9.8}{m.s^{-2}}$.}{\fgch{1}}
  \fgrow{\fgNo}{$t = \sqrt{\dfrac{2h}{g}}$ \quad $R = u\sqrt{\dfrac{2h}{g}}$}
    {Horizontal launch from a height $h$.}{\fgch{1}}
  \fgrow{\fgNo}{$t_\mathrm{top} = \dfrac{u_y}{g}$ \quad $h_\mathrm{max} = \dfrac{u_y^2}{2g}$}
    {Time to the top; greatest height.}{\fgch{1}}
  \fgrow{\fgNo}{$t_\mathrm{flight} = \dfrac{2u_y}{g}$}
    {Time of flight on level ground.}{\fgch{1}}
  \fgrow{\fgNo}{$R = \dfrac{u^2\sin 2\theta}{g}$}
    {Range on level ground (largest at \ang{45}).}{\fgch{1}}
  \fgrow{\fgNo}{$v = \sqrt{v_x^2 + v_y^2}$ \quad $\tan\alpha = \dfrac{\lvert v_y\rvert}{v_x}$}
    {Speed and direction at landing.}{\fgch{1}}
  \fgrow{\fgComb}{$-h = u_y t - \dfrac{1}{2}gt^2$}
    {Launch from a height $h$ (up positive): a quadratic in $t$; keep the positive root.}{\fgch{1}}
  \fgrow{\fgComb}{$\dfrac{1}{2}mv_2^2 - \dfrac{1}{2}mv_1^2 = mg(h_1 - h_2)$}
    {Speed at a known height (no direction).}{\fgch{1}}
  \fgrow{\fgComb}{$s = ut$ \quad $v = u$}
    {Uniform velocity ($a = 0$).}{\fgch{1}}
  \fgnext
  % ---- Chapter 2: Circular Motion: Forces
  \fgrow{\fgYes}{$v = \dfrac{2\pi r}{T}$}{Speed in a circle.}{\fgch{2}}
  \fgrow{\fgYes}{$a_\mathrm{c} = \dfrac{v^2}{r}$}{Acceleration towards the centre.}{\fgch{2}}
  \fgrow{\fgYes}{$F_\mathrm{c} = \dfrac{mv^2}{r}$}{Net force towards the centre (real forces
    make it).}{\fgch{2}}
  \fgrow{\fgNo}{$\tan\theta = \dfrac{v^2}{rg}$}{Conical pendulum; banked road at its design
    speed; bowl.}{\fgch{2}}
  \fgrow{\fgNo}{$v_\mathrm{max} = \sqrt{\mu g r}$}{Fastest speed on a flat bend.}{\fgch{2}}
  \fgrow{\fgNo}{$v_\mathrm{min} = \sqrt{\dfrac{gr}{\mu}}$}{Slowest safe speed in a rotor.}
    {\fgch{2}}
  \fgnext
  % ---- Chapter 3: Circular Motion: Angles, Energy, Work and Torque
  \fgrow{\fgYes}{$\omega = \dfrac{\Delta\theta}{t}$}{Angular velocity ($\Delta\theta$ in
    radians).}{\fgch{3}}
  \fgrow{\fgNo}{$\omega = \dfrac{2\pi}{T} = 2\pi f$}{Angular velocity from $T$ or $f$.}
    {\fgch{3}}
  \fgrow{\fgNo}{$v = \omega r$}{Speed of a point on a spinning object.}{\fgch{3}}
  \fgrow{\fgNo}{$a_\mathrm{c} = \omega^2 r$}{Centripetal acceleration of a point on a spinning
    object.}{\fgch{3}}
  \fgrow{\fgYes}{$W = F_\parallel s = Fs\cos\theta$}{Work by a force ($F_\mathrm{c}$ does
    none).}{\fgch{3}}
  \fgrow{\fgNo}{$W_\mathrm{other} = \Delta E = \Delta K + \Delta U$}{Work by forces other than
    gravity.}{\fgch{3}}
  \fgrow{\fgComb}{$\tfrac{1}{2}mv_1^2 + mgh_1 = \tfrac{1}{2}mv_2^2 + mgh_2$}{Speeds on a
    vertical circle (no friction).}{\fgch{3}}
  \fgrowii{\fgNo}{top: $F_\mathrm{T} + mg = \dfrac{mv^2}{r}$}
    {bottom: $F_\mathrm{T} - mg = \dfrac{mv^2}{r}$}{Tension (or the normal force inside a
    loop) at the very top or bottom of a vertical circle.}{\fgch{3}}
  \fgrowii{\fgNo}{dip: $F_\mathrm{N} - mg = \dfrac{mv^2}{r}$}
    {crest: $mg - F_\mathrm{N} = \dfrac{mv^2}{r}$}{Normal force (apparent weight) in a dip or
    on a crest.}{\fgch{3}}
  \fgrow{\fgYes}{$\tau = r_\perp F = rF\sin\theta$}{Torque ($\theta$ between $r$ and $F$).}
    {\fgch{3}}
  \fgrow{\fgNo}{$\sum\tau_\mathrm{clockwise} = \sum\tau_\mathrm{anticlockwise}$}{Balanced
    torques: the rotation does not change.}{\fgch{3}}
  \fgnext
  % ---- Chapter 4: Gravity, Orbits and Kepler's Laws
  \fgrow{\fgYes}{$F = \dfrac{GMm}{r^2}$}{Gravitational force ($r$ centre to centre).}
    {\fgch{4}}
  \fgrow{\fgNo}{$r = r_\mathrm{E} + h$}{Distance from the centre, from a height $h$.}{\fgch{4}}
  \fgrow{\fgNo}{$g = \dfrac{GM}{r^2}$}{Field strength at or above a planet.}{\fgch{4}}
  \fgrow{\fgNo}{$v = \sqrt{\dfrac{GM}{r}}$}{Speed in a circular orbit.}{\fgch{4}}
  \fgrow{\fgYes}{$\dfrac{r^3}{T^2} = \dfrac{GM}{4\pi^2}$}{Kepler's third law: find $T$, $r$ or
    $M$.}{\fgch{4}}
  \fgrow{\fgNo}{$r_\mathrm{p} v_\mathrm{p} = r_\mathrm{a} v_\mathrm{a}$}{Kepler's second law
    (closest, farthest point).}{\fgch{4}}
  \fgrow{\fgComb}{$r = \sqrt[3]{\dfrac{GMT^2}{4\pi^2}}$ \quad $h = r - r_\mathrm{E}$}
    {Orbit radius and height from the period.}{\fgch{4}}
  \fgnext
  % ---- Chapter 5: Energy in Gravitational Fields
  \fgrow{\fgYes}{$U = -\dfrac{GMm}{r}$}{Potential energy (zero at infinity).}{\fgch{5}}
  \fgrowii{\fgNo}{$K = \dfrac{GMm}{2r}$}{$E = U + K = -\dfrac{GMm}{2r}$}{Kinetic and total
    energy of a \emph{circular} orbit.}{\fgch{5}}
  \fgrow{\fgComb}{$E = K + U = \dfrac{1}{2}mv^2 - \dfrac{GMm}{r}$}{Total energy on any path
    with gravity only. $E < 0$: bound.}{\fgch{5}}
  \fgrow{\fgNo}{$v_\mathrm{esc} = \sqrt{\dfrac{2GM}{r}}$}{Escape velocity (no engine).}
    {\fgch{5}}
  \fgrow{\fgNo}{$\Delta E = \dfrac{GMm}{2}\left(\dfrac{1}{r_1} - \dfrac{1}{r_2}\right)$}
    {Energy to change circular orbit.}{\fgch{5}}
  \fgnext
  % ---- Year 11 formulas and rearranged forms
  \multicolumn{4}{@{}l@{}}{\rule[-1.1ex]{0pt}{4.2ex}\normalsize
    {\bfseries\color{hscbrand}Other formulas used in this module} (Year 11, or rearranged
    forms)}\\\hline
  \fgrow{\fgYes}{$\vec{F}_\mathrm{net} = m\vec{a}$}{Newton's second law.}{\fgch{1}, \fgch{2}}
  \fgrow{\fgYes}{$K = \dfrac{1}{2}mv^2$}{Kinetic energy.}{\fgch{1}, \fgch{3}}
  \fgrow{\fgYes}{$\Delta U = mg\Delta h$}{Change in $U$ near the surface only.}
    {\fgch{1}, \fgch{3}}
  \fgrow{\fgYes}{$f = \dfrac{1}{T}$}{Frequency from the period.}{\fgch{2}}
  \fgrow{\fgNo}{$F_\mathrm{f} \le \mu F_\mathrm{N}$}{Static friction is at most
    $\mu F_\mathrm{N}$.}{\fgch{2}}
  \fgrow{\fgComb}{$a_\mathrm{c} = \dfrac{4\pi^2 r}{T^2}$}{Centripetal acceleration from the
    period.}{\fgch{2}}
  \fgrow{\fgNo}{$s = r\theta$}{Arc length ($\theta$ in radians).}{\fgch{3}}
  \fgrow{\fgNo}{$W_\mathrm{net} = \Delta K$}{Work--energy theorem.}{\fgch{3}}
  \fgrow{\fgNo}{$v = \sqrt{gr}$}{Smallest speed at the top (string slack).}{\fgch{3}}
  \fgrow{\fgComb}{$T = 2\pi\sqrt{\dfrac{r^3}{GM}}$}{Period of a circular orbit.}{\fgch{4}}
  \fgrow{\fgComb}{$\dfrac{r_1^3}{T_1^2} = \dfrac{r_2^3}{T_2^2}$}{Two orbits around the
    \emph{same} central body.}{\fgch{4}}
  \fgrow{\fgComb}{$\Delta U = GMm\left(\dfrac{1}{r_1} - \dfrac{1}{r_2}\right)$}{Change in $U$
    from $r_1$ to $r_2$ ($U_2 - U_1$).}{\fgch{5}}
  \fgrow{\fgNo}{$W_\text{by gravity} = -\Delta U$}{Gravity's work and the change in $U$.}
    {\fgch{3}, \fgch{5}}
\end{longtable}}

\section*{Data sheet values used in this module}
The exam gives you the data sheet. These are the values that Module~\modulenumber\ uses,
exactly as printed.

\begin{table}
  \centering
  \small
  \begin{tabular}{@{}L{58mm}L{50mm}L{32mm}@{}}
    \toprule
    \colhead{Quantity (as on the sheet)} & \colhead{Value} & \colhead{Used in}\\
    \midrule
    Earth's gravitational acceleration & $g = \qty{9.8}{m.s^{-2}}$ & \cref{ch1:chap,ch2:chap,ch3:chap,ch4:chap,ch5:chap}\\
    Universal gravitational constant & $G = \qty{6.67e-11}{N.m^2.kg^{-2}}$ & \cref{ch4:chap,ch5:chap}\\
    Mass of Earth & $M_\mathrm{E} = \qty{6.0e24}{kg}$ & \cref{ch4:chap,ch5:chap}\\
    Radius of Earth & $r_\mathrm{E} = \qty{6.371e6}{m}$ & \cref{ch4:chap,ch5:chap}\\
    Speed of light & $c = \qty{3.00e8}{m.s^{-1}}$ & \cref{ch4:chap} (signal delay)\\
    \bottomrule
  \end{tabular}
\end{table}

\begin{itemize}
  \item The masses and radii of the Sun, the Moon and other planets are \emph{not} on the data
        sheet. When they are needed, they are given to you.
  \item $GM_\mathrm{E}/r_\mathrm{E}^2$ gives $g = \qty{9.86}{m.s^{-2}}$, because $M_\mathrm{E}$ is
        rounded on the sheet. Use \qty{9.8}{m.s^{-2}} unless you must calculate $g$.
  \item Use SI units: \si{km} $\times 1000$ gives \si{m}; \si{km.h^{-1}} $\div 3.6$ gives
        \si{m.s^{-1}}; hours $\times 3600$ and days $\times \num{86400}$ give \si{s};
        degrees $\times \pi/180$ give \si{rad}.
\end{itemize}
}{}
\IfFileExists{./experiments.tex}{% experiments.tex -- Module 5 ESSENTIALS: Key experiments at a glance.
% One row per key experiment of the module. Facts and numbers are those of the Key
% experiment boxes in ch1.tex ... ch5.tex (and of the full book, src/mod5/).
% Every row has a Key experiment box in its chapter (the orbit simulation: ch4, Section 4.5,
% with the full book's numbers).
\unnumberedchapter[\faFlask]{Key experiments at a glance}
\phantomsection\label{kx:chap}

% ---- helpers for this part only (prefix kx)
\newcommand\kxhead{\colhead{Experiment} & \colhead{What was done} & \colhead{What it showed} &
  \colhead{Ch.}}
% \kxrow{experiment}{dot point}{what was done}{what it showed}{chapter number}
\newcommand\kxrow[5]{\textbf{#1}\newline{\footnotesize\color{hscmuted}#2} & #3 & #4 &
  \hyperref[ch#5:chap]{\textbf{#5}}\\\hline}

The syllabus asks you to carry out or study these investigations. Each one has a Key
experiment box in its chapter, with its aim, method, result, conclusion and significance.
\textbf{Required} marks the two practical investigations that the syllabus names.

{\small
\setlength\LTpre{4pt}\setlength\LTpost{6pt}
\renewcommand\arraystretch{1.25}\arrayrulecolor{hscrule}
\begin{longtable}{@{}L{35mm}@{\hspace{3mm}}L{48mm}@{\hspace{3mm}}%
    L{\dimexpr\linewidth-35mm-48mm-7mm-9mm\relax}@{\hspace{3mm}}C{7mm}@{}}
  \kxhead\\\hline\endfirsthead
  \kxhead\\\hline\endhead
  \kxrow{Range of a ball launched horizontally (required)}{M5.IQ1.03}
    {A steel ball leaves a bench ramp horizontally at 5 speeds ($u$ from a light gate); the
     range $R$ is measured.}
    {$R$ against $u$ is a straight line very close to the origin, gradient \qty{0.423}{s} (model
     \qty{0.429}{s}). This validates $R = u\sqrt{2h/g}$: the two motions are independent.}{1}
  \kxrow{Whirling a stopper on a string (required)}{M5.IQ2.01}
    {A stopper is whirled through a tube; a hanging mass gives $F_\mathrm{c} = Mg$; 20 turns are
     timed. Change $M$, then $r$, then $m$.}
    {Part A: $F_\mathrm{c} \propto v^2$ (a straight line through the origin). Parts B and C
     (expected): $F_\mathrm{c} \propto 1/r$ and $F_\mathrm{c} \propto m$.}{2}
  \kxrow{Torque on a metre rule}{M5.IQ2.05}
    {A pivoted metre rule is held level by a spring balance at different distances and
     angles; then the torques are unbalanced.}
    {$Fr$ stays the same (about \qty{0.59}{N.m}); only $F\sin\theta$ turns it:
     $\tau = rF\sin\theta$. A bigger net torque changes the rotation faster.}{3}
  \kxrow{Work and energy in circular motion}{M5.IQ2.04}
    {A whirled stopper is filmed; a ball swinging down from rest is timed with a light gate.}
    {The tension does no work: $K$ is constant in a horizontal circle. In the swing, $K$
     gained $=$ $U$ lost (within a 3\,\% light-gate error).}{3}
  \kxrow{Factors that affect the field strength $g$}{M5.IQ3.01}
    {Secondary data: surface $g$ against $M/r^2$ for the Moon, Mercury, Mars, Venus and the
     Earth.}
    {A straight line through the origin with gradient $G$, so $g = GM/r^2$.}{4}
  \kxrow{Kepler's third law and the moons of Jupiter}{M5.IQ3.04}
    {Secondary data: $r^3$ against $T^2$ for Io, Europa, Ganymede and Callisto.}
    {One value of $r^3/T^2$ for all four moons. The gradient, $GM/4\pi^2$, gives Jupiter's
     mass, \qty{1.90e27}{kg}.}{4}
  \kxrow{Orbits in a simulation}{M5.IQ3.02}
    {A computer orbit model: change only the satellite's mass, then only the orbital radius.}
    {The period does not depend on the satellite's mass. Bigger orbits are slower:
     $v \propto 1/\sqrt{r}$.}{4}
  \kxrow{Launch speed, total energy and escape (simulation)}{M5.IQ3.05}
    {A probe is launched sideways from $r_0 = \qty{1.00e7}{m}$ at 5 to \qty{9}{km.s^{-1}} in a
     computer orbit model.}
    {It falls back, orbits or escapes. It escapes only when $E \ge 0$: between 8.94 and
     \qty{8.95}{km.s^{-1}}, and $\sqrt{2GM/r_0} = \qty{8.95}{km.s^{-1}}$.}{5}
\end{longtable}}
}{}
\IfFileExists{./glossary.tex}{% glossary.tex -- Module 5 ESSENTIALS glossary (A-Z): the key terms only.
% src/mod5/glossary.tex (term, meaning and Chinese copied word for word, except the
% rows corrected in the script's FIXES). Edit the script, not this file.
\unnumberedchapter[\faBook]{Glossary}

The key words of this module, from A to Z. The meaning is in plain English. The last column
gives the usual Chinese term used in physics textbooks. The full guide's glossary has more
words.

\begin{glossarytable}
  % ---- A
  \gterm{Acceleration due to gravity ($g$)}{the acceleration of a falling object near the Earth: \qty{9.8}{m.s^{-2}} downwards}{重力加速度}
  \gterm{Angular displacement ($\Delta\theta$)}{the angle an object turns through, measured at the centre}{角位移}
  \gterm{Angular velocity ($\omega$)}{the angle turned each second, in \si{rad.s^{-1}}}{角速度}
  \gterm{Aphelion}{the farthest point of an orbit from the Sun}{远日点}
  \gterm{Apogee}{the farthest point of an orbit around the Earth}{远地点}
  \gterm{Apparent weightlessness}{the feeling of having no weight in free fall, because the normal force is zero}{完全失重}
  % ---- B
  \gterm{Balanced torques}{clockwise and anticlockwise torques are equal, so the rotation does not change}{力矩平衡}
  \gterm{Banked track}{a curved road that is tilted, higher on the outside of the bend}{倾斜弯道（外高内低）}
  \gterm{Bound}{trapped by gravity, so the object cannot get away; its total energy is negative}{被引力束缚的}
  % ---- C
  \gterm{Centrifugal force}{a made-up ``outward force''; do not use it: the outward feeling is inertia}{离心力}
  \gterm{Centripetal acceleration ($a_\mathrm{c}$)}{the acceleration towards the centre of a circle: $v^2/r$}{向心加速度}
  \gterm{Centripetal force ($F_\mathrm{c}$)}{the net force towards the centre that keeps an object moving in a circle}{向心力}
  \gterm{Coefficient of friction ($\mu$)}{a number with no unit that says how strongly two surfaces grip each other}{摩擦因数}
  \gterm{Component}{the part of a vector (such as a force) that points in one chosen direction}{分量（分力）}
  \gterm{Conical pendulum}{a mass on a string moving in a horizontal circle; the string sweeps out a cone}{圆锥摆}
  % ---- E
  \gterm{Ellipse}{a stretched circle (an oval)}{椭圆}
  \gterm{Escape velocity ($v_\mathrm{esc}$)}{the smallest launch speed that lets a projectile get away from a planet for ever}{逃逸速度}
  % ---- F
  \gterm{Free fall}{motion when gravity is the only force acting}{自由下落（只受引力）}
  % ---- G
  \gterm{Geostationary orbit}{a circular orbit above the equator with a one-day period, so the satellite stays above one point}{地球静止轨道}
  \gterm{Geosynchronous orbit}{any orbit with a period of one day; it may be tilted}{地球同步轨道}
  \gterm{Gravitational field strength ($g$)}{the gravitational force on each kilogram at a point, in \si{N.kg^{-1}} (\si{m.s^{-2}})}{引力场强度}
  \gterm{Gravitational force}{the pull between any two objects that have mass}{引力（万有引力）}
  \gterm{Gravitational potential energy ($U$)}{energy stored because of where an object is in a gravitational field; near the ground it changes by $mg\Delta h$; in space $U = -GMm/r$}{引力势能（地面附近：重力势能）}
  % ---- H
  \gterm{Hohmann transfer orbit}{a fuel-saving half-ellipse between two circular orbits, using two short engine burns (both forwards to move up, both backwards to move down)}{霍曼转移轨道}
  % ---- I
  \gterm{Inertia}{the tendency of an object to stay at rest, or to keep moving in a straight line at a steady speed}{惯性}
  \gterm{Infinity}{a place so far away that the planet's gravity there is too small to matter}{无穷远处}
  \gterm{Inverse-square law}{a rule where doubling the distance makes the effect four times smaller}{平方反比定律}
  % ---- K
  \gterm{Kepler's first law (law of ellipses)}{each planet moves in an ellipse with the Sun at one focus}{开普勒第一定律（轨道定律）}
  \gterm{Kepler's second law (law of equal areas)}{the line from the Sun to a planet sweeps out equal areas in equal times}{开普勒第二定律（面积定律）}
  \gterm{Kepler's third law (law of periods)}{$T^2$ is proportional to $r^3$ for everything that orbits the same body}{开普勒第三定律（周期定律）}
  \gterm{Kinetic energy ($K$)}{energy of motion: $K = \tfrac{1}{2}mv^2$}{动能}
  % ---- L
  \gterm{Line of action}{the line along which a force acts, drawn longer in both directions}{力的作用线}
  \gterm{Low Earth orbit (LEO)}{an orbit about \qtyrange{200}{2000}{km} above the Earth, with a period of about 90 minutes to 2 hours}{低地球轨道（近地轨道）}
  % ---- M
  \gterm{Maximum height ($h_\mathrm{max}$)}{the greatest height above the launch point}{最大高度}
  % ---- N
  \gterm{Net force}{the total of all the forces on an object, added with their directions}{合力}
  \gterm{Normal force ($F_\mathrm{N}$)}{the push of a surface on an object, at \ang{90} to the surface}{支持力（法向力）}
  % ---- O
  \gterm{Orbital decay}{the slow fall to lower orbits caused by drag from the thin upper atmosphere}{轨道衰减}
  \gterm{Orbital speed ($v$)}{the speed a satellite needs to stay in a circular orbit}{环绕速度（轨道速度）}
  % ---- P
  \gterm{Perigee}{the closest point of an orbit around the Earth}{近地点}
  \gterm{Perihelion}{the closest point of an orbit to the Sun}{近日点}
  \gterm{Period ($T$)}{the time for one full turn}{周期}
  \gterm{Perpendicular distance ($r_\perp$)}{the shortest distance from the pivot to the line of action of a force (the lever arm)}{力臂}
  \gterm{Pivot}{the fixed point or line that an object turns about, or that a string hangs from}{支点（转轴；悬点）}
  \gterm{Projectile}{an object that is thrown or launched and then moves with no engine; in the model, gravity is the only force on it}{抛体}
  % ---- R
  \gterm{Radial field}{a field whose lines all point straight at the centre of the planet, like the spokes of a wheel}{径向场}
  \gterm{Radian (rad)}{a unit for angles: one full turn is $2\pi$ rad, the same as \ang{360}}{弧度}
  \gterm{Range ($R$)}{the horizontal distance from the launch point to the landing point}{射程（水平射程）}
  % ---- S
  \gterm{Semi-major axis ($a$)}{half of the longest diameter of an ellipse}{半长轴}
  \gterm{Static friction ($F_\mathrm{f}$)}{friction that stops surfaces sliding; it grows as needed, up to a maximum of $\mu F_\mathrm{N}$}{静摩擦力}
  % ---- T
  \gterm{Tangent}{a line that touches a circle at one point, at \ang{90} to the radius}{切线}
  \gterm{Tension ($F_\mathrm{T}$)}{the pull of a string, rope or chain, along its length}{拉力（张力）}
  \gterm{Time of flight ($t_\mathrm{flight}$)}{the time from launch until landing}{飞行时间}
  \gterm{Torque ($\tau$)}{the turning effect of a force, in \si{N.m}}{力矩}
  \gterm{Total mechanical energy ($E$)}{kinetic energy plus potential energy: $E = K + U$}{机械能}
  % ---- U
  \gterm{Uniform circular motion}{moving around a circle at a constant speed}{匀速圆周运动}
  \gterm{Uniform velocity}{a velocity that does not change: the same speed in the same direction; the acceleration is zero}{匀速直线运动（速度恒定）}
  \gterm{Universal gravitational constant ($G$)}{the fixed number that sets how strong gravity is, \qty{6.67e-11}{N.m^2.kg^{-2}}}{引力常量}
  % ---- W
  \gterm{Work ($W$)}{energy transferred by a force when the object moves along the direction of the force (or part of it)}{功}
  \gterm{Work--energy theorem}{the work done by the net force equals the change in kinetic energy}{动能定理}
\end{glossarytable}
}{}

\end{document}
